\documentclass[11pt,a4paper]{article}
\usepackage{amsmath,amssymb,amsthm}
\usepackage{geometry}
\usepackage{hyperref}
\usepackage{cascade-xref}        % cross-paper \xref machinery (incl. xr-hyper)
\usepackage{booktabs}
\xrhyperdoc{part4b}{cascade-series-part4b}
\geometry{margin=2.5cm}

\newtheorem{theorem}{Theorem}[section]
\newtheorem{lemma}[theorem]{Lemma}
\newtheorem{corollary}[theorem]{Corollary}
\newtheorem{proposition}[theorem]{Proposition}
\newtheorem{definition}[theorem]{Definition}
\newtheorem{remark}[theorem]{Remark}

\title{The Cascade Series --- Part 0\\[0.5em]
\large Scale Variance from Orthogonality:\\
How the Unit Ball Generates $10^{120}$ Orders of Magnitude}
\author{RTAC}
\date{March 2026}

\begin{document}
\setlength{\emergencystretch}{3em}
\maketitle
\phantomsection
\label{paperdoc}

\begin{abstract}
The unit ball in $\mathbb{R}^d$ has a slicing recurrence
$V_{d+1}/V_d = \sqrt{\pi}\cdot R(d+1)$ containing one constant,
$\sqrt{\pi} = \Gamma(1/2)$, forced by orthogonality. The sphere-area decay
rate $p(d)$ has a unique natural zero at $d_0 = 7$. This zero generates
two threshold values $c_1 = \tfrac{1}{2}\ln\pi$ and $c_2 = \sqrt{\pi}$ ---
the two canonical values of the recurrence's unique constant, one in each
of its two canonical forms. The orbit terminates because a first-order
recurrence has exactly two canonical forms and one bridge map between them.
The thresholds select $d_1 = 19$ and $d_2 = 217$.

The three privileged values of $p(d)$---zero, $c_1$, and $c_2$---partition
the cascade into exactly four regimes (Growth, Decay, Exponential decay,
Oblivion), separated by three continuous boundary crossings lying strictly
between integers: $d_0^\star\approx 6.257$, $d_1^\star\approx 19.731$,
$d_2^\star\approx 217.627$. The Gamma function thus produces exactly four
distinguished dimensions: the volume maximum $d_V = 5$ (an interior
landmark of the Growth regime), plus one representative integer per
boundary crossing. No fifth exists.

The four dimensions separate by structural stability into two classes:
the equilibria are structurally stable (they exist for any value of the
recurrence constant) and serve as scale; the thresholds are structurally
sensitive (they depend on the specific constant $\sqrt{\pi}$) and carry
content. Scale enters as a ratio; content enters as a product.

The cascade invariant
\[
I = \max_{(d_0,d_1,d_2)}\;\frac{\Omega_5^2}{\Omega_{d_0}^2}\;\Omega_{d_1}\;\Omega_{d_2}
= \frac{9}{\pi^2}\;\Omega_{19}\;\Omega_{217}
= 1.0989\times 10^{-120},
\]
where the supremum runs over integer labelings of the three boundary pairs
$(d_0,d_1,d_2) \in \{6,7\}\times\{19,20\}\times\{217,218\}$ and the unique
argmax is $(7, 19, 217)$. This variational definition eliminates the need
for an integer rounding convention at each boundary pair. $I$ is the
unique order-0 (global) cascade invariant; it depends only on the four
distinguished dimensions and is immune to subleading corrections from
inter-layer structure. The formula for the hierarchy
is $\pi e^{2\sqrt{\pi}}(2\sqrt{\pi}-1)/\ln 10\approx 120$.

Every step is a theorem about the Gamma function. No physics enters.
\end{abstract}

\medskip\noindent\textbf{Notation.}
Throughout this paper, $\Omega_d$ denotes the area of $S^d$:
$\Omega_d = 2\pi^{(d+1)/2}/\Gamma((d+1)/2)$. The boundary of the unit
ball $B^d$ is $S^{d-1}$ with area $\Omega_{d-1}$. Boundary dominance
(Theorem~\ref{thm:boundary}): $\Omega_{d-1}/V_d = d$.

% ============================================================
\section{The Question}\label{sec:the-question}
% ============================================================

How much scale variance can pure mathematics produce from nothing?

The unit ball $B^d = \{x\in\mathbb{R}^d : |x|\leq 1\}$ exists in every
dimension. The area of the unit sphere $S^d$ is
$\Omega_d = 2\pi^{(d+1)/2}/\Gamma((d+1)/2)$. At $d = 7$:
$\Omega_7 = \pi^4/3\approx 32.5$. At $d = 217$:
$\Omega_{217}\approx 10^{-120}$. The ratio exceeds $10^{121}$. The unit
ball generates 121 orders of magnitude of scale variance with no free
parameters and no input beyond the definition of orthogonal dimensions.

This paper extracts two distinguished scales from this variance, proves
they combine uniquely with the cascade's two natural equilibria, and evaluates
the result.

\medskip
\noindent\emph{Read upstream from the Prelude.} The Prelude argues that
the slicing recurrence developed below \emph{is} the resolution mechanism
by which the asymptotic upper edge $B^\infty$ extracts finite content at
each integer $d$, and that the recurrence and its tower are forced by
orthogonality plus $\mathrm{RCA}_0$-level minimality. This paper develops
what the recurrence produces. The arguments below are theorems about the
Gamma function and stand independently of the Prelude's foundational
chain.

% ============================================================
\section{The Recurrence}\label{sec:recurrence}
% ============================================================

Cut $B^{d+1}$ perpendicular to one axis. Each cross-section at height $x$
is a $d$-ball of radius $\sqrt{1-x^2}$:
\begin{equation}\label{eq:slice}
V_{d+1} = V_d\int_{-1}^{1}(1-x^2)^{d/2}\,dx.
\end{equation}
Substitute $x = \cos\theta$. The integral becomes
$2\int_0^{\pi/2}\sin^d\theta\,d\theta$: one quarter turn from axis to
equator. Via the Beta function:
\begin{equation}\label{eq:recurrence}
V_{d+1} = V_d\cdot\sqrt{\pi}\cdot R(d+1), \qquad
R(d) := \frac{\Gamma((d+1)/2)}{\Gamma((d+2)/2)}.
\end{equation}

\begin{theorem}[Orthogonal compression constant]\label{thm:sqrtpi}
$\Gamma(1/2) = \sqrt{\pi}$ is the unique dimension-independent constant in
the recurrence, forced by orthogonality. It appears with coefficient
exactly~1.
\end{theorem}

\begin{proof}
The first argument of $B$ is $\frac{1}{2}$, forced by orthogonality: the
angle between axis and equator is $\pi/2$; the substitution
$x = \cos\theta$ produces the half-power $t^{-1/2}$ in $B(1/2,\cdot)$;
$\Gamma(1/2) = \sqrt{\pi}$. The chain: dimension $\to$ orthogonal $\to$
$\pi/2$ $\to$ $B(1/2,\cdot)$ $\to$ $\sqrt{\pi}$. Every arrow forced.
\end{proof}

\begin{lemma}[Slicing-ratio duplication]\label{lem:R-duplication}
For every integer $d \geq 1$,
\[
R(2d-1)\cdot R(2d) \;=\; \frac{1}{d}.
\]
Equivalently, the cascade-tower product of slicing ratios telescopes to a factorial:
\[
\prod_{d=1}^{2N} R(d) \;=\; \frac{1}{N!}.
\]
\end{lemma}

\begin{proof}
By the definition $R(d) = \Gamma((d+1)/2)/\Gamma((d+2)/2)$,
\[
R(2d-1)\;=\;\frac{\Gamma(d)}{\Gamma(d+\tfrac{1}{2})},\qquad
R(2d)\;=\;\frac{\Gamma(d+\tfrac{1}{2})}{\Gamma(d+1)}.
\]
Their product telescopes through the half-integer Gamma factor: $R(2d-1)R(2d) = \Gamma(d)/\Gamma(d+1) = 1/d$. The product over $d = 1, \ldots, N$ gives $\prod_{d=1}^{2N} R(d) = \prod_{d=1}^{N} \frac{1}{d} = \frac{1}{N!}$. (This is the Legendre duplication formula $\Gamma(2z) = (2\pi)^{-1/2}\,2^{2z-1/2}\,\Gamma(z)\Gamma(z+\tfrac{1}{2})$ specialised to integer $z = d$, expressed in cascade slicing ratios.)
\end{proof}

\begin{remark}
The identity is structural, not approximate: cascade slicing ratios at consecutive odd-even layer pairs $(2d{-}1, 2d)$ satisfy a fixed product law $1/d$. The full-tower product $\prod_{d=1}^{2N} R(d) = 1/N!$ records this telescoping. Verified numerically in \texttt{tools/research/legendre\_duplication.py}.
\end{remark}

% ============================================================
\section{Sphere Areas Are Primary}\label{sec:sphere-areas-are-primary}
% ============================================================

\begin{theorem}[Boundary dominance]\label{thm:boundary}
$\Omega_{d-1}/V_d = d$ for all $d\geq 1$.
\end{theorem}

\begin{proof}
$\Omega_{d-1} = 2\pi^{d/2}/\Gamma(d/2)$,\;
$V_d = \pi^{d/2}/[(d/2)\Gamma(d/2)]$.\; Ratio: $d$.
\end{proof}

\begin{corollary}\label{cor:primary}
$V_d = \Omega_{d-1}/d$. Volumes are derived from sphere areas. The
recurrence is a first-order multiplicative map between sphere areas:
\[
\Omega_{d+1} = \Omega_d\cdot\sqrt{\pi}\cdot R(d).
\]
Sphere area is the unique independent cascade quantity at each level.
Every other cascade object---$V_d$, $R(d)$,
$R_{\mathrm{eff}}(d) = 1/\sqrt{d+3}$ (Theorem~\ref{thm:reff})---is
derived from sphere areas.
\end{corollary}

% ============================================================
\section{Irreducibility}\label{sec:irreducibility}
% ============================================================

\begin{theorem}\label{thm:irred}
$\sqrt{\pi}$ is the irreducible residue of the recurrence.
\end{theorem}

\begin{proof}
The cascade's first-order multiplicative recurrence on sphere areas
(Corollary~\ref{cor:primary}) is
\[
\frac{\Omega_{d+1}}{\Omega_d} \;=\; \sqrt{\pi}\cdot R(d),
\qquad R(d) := \frac{\Gamma((d+1)/2)}{\Gamma((d+2)/2)}.
\]
This factors uniquely as a dimension-independent constant
$\sqrt{\pi}$ times a dimension-dependent Gamma ratio $R(d)$. The
logarithmic form
\[
\ln(\Omega_{d+1}/\Omega_d) \;=\; \tfrac{1}{2}\ln\pi + \ln R(d)
\]
isolates $\tfrac{1}{2}\ln\pi$ as the unique constant additive
component; the second log-difference
\[
\delta(d) \;:=\;
\ln(\Omega_{d+1}/\Omega_d) - \ln(\Omega_d/\Omega_{d-1})
\;=\; \ln R(d) - \ln R(d-1)
\]
removes the constant, leaving only Gamma-ratio content (the
$\tfrac{1}{2}\ln\pi$ term cancels at second order; $\pi$ persists
inside $R$ at half-integer arguments). Thus $\sqrt{\pi}$ is the
unique dimension-independent multiplicative content of the
recurrence; every other recurrence content lies in the
dimension-dependent Gamma ratios.
\end{proof}

\begin{theorem}[Effective slicing radius]\label{thm:reff}
The root-mean-square height of the slicing measure $(1-x^2)^{d/2}$ on
$[-1,1]$ is
\[
R_{\mathrm{eff}}(d) := \sqrt{\langle x^2\rangle_d} = \frac{1}{\sqrt{d+3}},
\qquad
\langle x^2\rangle_d := \frac{\int_{-1}^{1} x^2(1-x^2)^{d/2}\,dx}
                              {\int_{-1}^{1}(1-x^2)^{d/2}\,dx}.
\]
\end{theorem}

\begin{proof}
Both integrals are Beta functions:
$\int_{-1}^{1}(1-x^2)^{d/2}\,dx = B(\tfrac{1}{2},\tfrac{d}{2}+1)$ and
$\int_{-1}^{1}x^2(1-x^2)^{d/2}\,dx = B(\tfrac{3}{2},\tfrac{d}{2}+1)$
(substitute $t = x^2$ in each, as in \eqref{eq:slice}). Their ratio is
\[
\frac{B(\tfrac{3}{2},\tfrac{d}{2}+1)}{B(\tfrac{1}{2},\tfrac{d}{2}+1)}
= \frac{\Gamma(\tfrac{3}{2})}{\Gamma(\tfrac{1}{2})}\cdot
  \frac{\Gamma(\tfrac{d}{2}+\tfrac{3}{2})}{\Gamma(\tfrac{d}{2}+\tfrac{5}{2})}
= \frac{1}{2}\cdot\frac{2}{d+3}
= \frac{1}{d+3}.
\]
Taking the square root gives $R_{\mathrm{eff}}(d) = 1/\sqrt{d+3}$.
\end{proof}

\begin{remark}
$R_{\mathrm{eff}}(d) > 0$ for every finite $d$: the slicing concentrates
near the equator but never collapses to it. This is the per-step
compactification residual---the scale at which consecutive slicing
layers fail to be perfectly collinear in $L^2$. \S\ref{sec:gram-closed-form}
quantifies this overlap directly: the adjacent-layer overlap deficit
scales as $1 - C^2_{d,d+1}\sim 1/(8d^2) = R_{\mathrm{eff}}^4/8 +
O(1/d^3)$ (Theorem~\ref{thm:overlap}).
\end{remark}

% ============================================================
\section{The Natural Zero}\label{sec:the-natural-zero}
% ============================================================

The sphere-area decay rate decomposes as:
\[
p(d) = -\tfrac{1}{2}\ln\pi +
\tfrac{1}{2}\psi\!\left(\tfrac{d+1}{2}\right),
\]
where $\psi$ is the digamma function. Since $\psi^{(1)} > 0$, $p$ is
strictly monotone increasing with a unique zero.

\begin{definition}[Natural zero]\label{def:d0}
$d_0$ is the smallest integer $d$ such that $p(d) > 0$: the first
integer in the decay regime. Stirling: the continuous zero of $p$ lies at
$d^*\approx 2\pi\approx 6.27$; integer $d_0 = 7$.
\end{definition}

\begin{theorem}[Dual characterisation of $d_0$]\label{thm:dual}
The natural zero $d_0$ has two independent characterisations:
\begin{enumerate}
\item[(a)] \emph{Decay onset.} $d_0$ is the first integer at which
$p(d) > 0$.
\item[(b)] \emph{Maximum-boundary ball.} $B^{d_0}$ is the unique ball
whose boundary sphere has maximum area among all unit spheres.
\end{enumerate}
Both give $d_0 = 7$.
\end{theorem}

\begin{proof}
Let $d_{\max} = \operatorname{argmax}\Omega_d$ (discrete). By
unimodality of $\Omega_d$, $d_{\max}$ is unique, and the discrete ratio
satisfies
\[
\Omega_{d_{\max}+1}/\Omega_{d_{\max}} < 1 < \Omega_{d_{\max}}/\Omega_{d_{\max}-1}.
\]
Since $p$ is strictly increasing and continuous, and its sign governs
the growth/decay transition of the continuous interpolation of $\Omega_d$,
the unique zero $d^*$ of $p$ lies in the interval
$(d_{\max},\;d_{\max}+1)$. Therefore the first integer with $p(d) > 0$ is
$d_0 = d_{\max} + 1$.

By boundary dominance (Theorem~\ref{thm:boundary}),
$\partial B^d = S^{d-1}$. Therefore
$\partial B^{d_0} = S^{d_{\max}}$, which carries the maximum sphere area.

Evaluation: $\Omega_6 = 16\pi^3/15\approx 33.1 > \Omega_7 = \pi^4/3
\approx 32.5 > \Omega_5 = \pi^3\approx 31.0$, so $d_{\max} = 6$.
Bracket: $p(6) = -0.021 < 0 < 0.056 = p(7)$. Both characterisations
give $d_0 = 7$.
\end{proof}

\begin{remark}\label{rem:convention}
The two characterisations are independent: (a)~uses only the
monotonicity of $p$; (b)~uses only the unimodality of $\Omega_d$ and
boundary dominance. Their agreement at $d_0 = 7$ is a consistency check
on the cascade structure, not a convention choice. The invariant
(Section~8) uses $\Omega_7$ because $p$ is defined at sphere dimensions,
its zero selects sphere dimension~7, and the invariant collects sphere
areas at the dimensions where $p$ takes distinguished values.
\end{remark}

% ============================================================
\section{The Threshold Pair}\label{sec:the-threshold-pair}
% ============================================================

\subsection{Operation A: the zero-crossing forces \texorpdfstring{$c_1$}{c\_1}}

\begin{theorem}[Operation A]\label{thm:opA}
The zero-crossing forces $c_1 := \tfrac{1}{2}\ln\pi$.
\end{theorem}

\begin{proof}
$p(d_0) = 0$ gives $\frac{1}{2}\psi((d_0+1)/2) = \frac{1}{2}\ln\pi$.
The right-hand side is the constant part of $p$, forced by
Theorem~\ref{thm:sqrtpi}.
\end{proof}

\begin{theorem}[Uniqueness of $c_1$]\label{thm:c1unique}
Among all scalars derivable from $d_0$, $c_1 = \tfrac{1}{2}\ln\pi$ is the
unique scalar satisfying (a) $c > 0$ (so $p(d) = c$ has a solution beyond
$d_0$) and (b) primitivity ($c$ is a recurrence constant).
\end{theorem}

\begin{proof}
\emph{Cascade-internal scalars at $d_0$ have the form $F(\sqrt{\pi}, d_0)$.}
By Theorem~\ref{thm:irred}, the slicing recurrence has a unique
nontrivial constant, $\sqrt{\pi}$. By Corollary~\ref{cor:primary}, sphere
area is the unique independent cascade quantity at each level, with
explicit closed form $\Omega_d = 2\pi^{(d+1)/2}/\Gamma((d+1)/2)$ --- an
analytic function of $\sqrt{\pi}$ and the integer~$d$. Every derived
cascade quantity inherits this dependence:
the volume $V_d = \Omega_{d-1}/d$ (Theorem~\ref{thm:boundary});
the Beta ratio $R(d) = \Gamma((d+1)/2)/\Gamma((d+2)/2)$;
the decay rate $p(d) = -\tfrac{1}{2}\ln\pi + \tfrac{1}{2}\psi((d+1)/2)$
together with its polygamma derivatives $p^{(n)}(d)$;
the effective slicing radius $R_{\mathrm{eff}}(d) = 1/\sqrt{d+3}$
(Theorem~\ref{thm:reff}). Each is an analytic function of $\sqrt{\pi}$
and~$d$. Sphere areas and derived quantities at integer-offset
dimensions $d_0 + k$ ($k\in\mathbb{Z}$) are analytic functions of
$\sqrt{\pi}$ and $d_0$ obtained by integer arithmetic on the dimension
label. Integer arithmetic on $d_0$ is structural rather than imported:
the slicing recurrence advances the dimension by $\pm 1$ per
application, so the orbit of $d_0$ under iteration of the recurrence
is exactly the integer lattice $\{d_0 + k : k\in\mathbb{Z}\}$, and any
$\Omega_{d_0+k}$ is reachable from $\Omega_{d_0}$ by $|k|$ recurrence
steps. Polynomial expressions in $d_0$ ($d_0+k$, $2d_0$, $d_0^2$,
etc.) are integer labels obtained by composing this stepping with the
data $d_0$ itself; the dimension label's arithmetic is the
recurrence's own discrete time, not an additional admissible
operation. The algebra of cascade-internal scalars derivable at $d_0$
is the closure of these primitives under analytic combination, hence
consists of analytic functions $F(\sqrt{\pi}, d_0)$.

\emph{Primitivity excludes $d_0$-dependence.} Criterion~(b) (``$c$ is a
recurrence constant'') is, by Definition~\ref{def:canonical} below,
the requirement that $c$ appear as the dimension-independent
multiplicative factor in a canonical form of the first-order
recurrence. Equivalently: $c$ does not depend on the integer label
$d$. A scalar $s = F(\sqrt{\pi}, d_0)$ for which $F$ depends nontrivially
on $d_0$ therefore fails~(b) by definition, even if $F$ evaluates to a
$d_0$-free numerical value at the specific dimension~$d_0$ (the same
expression at a different dimension would give a different value, so
the dimension is load-bearing in the construction). Primitive
candidates are restricted to functions $F$ independent of $d_0$ ---
i.e., to the sub-algebra of analytic functions $F(\sqrt{\pi})$ alone.

\emph{The four classes are an exhaustive source enumeration.} Within
the algebra of cascade-internal scalars at $d_0$, candidate primitives
arise from four natural sources, distinguished by which atomic input
the candidate is built from:
(i)~values of $p$ and its constant--variable decomposition at $d_0$;
(ii)~higher derivatives $p^{(n)}(d_0)$ for $n\geq 1$;
(iii)~the dimension label $d_0$ itself;
(iv)~functions of the primitive $\sqrt{\pi}$ alone (equivalently
$c_1 = \tfrac{1}{2}\ln\pi$ in log-space).
Sources~(i)--(iii) all introduce $d_0$-dependence into the candidate
expression and are excluded by the primitivity argument above; we
verify each below as a redundancy check, since some sources produce
$d_0$-independent \emph{values} despite being built from $d_0$-dependent
\emph{expressions} (e.g.\ $q(d_0) = \tfrac{1}{2}\ln\pi$ in
Class~1) and the verification identifies which of these values
coincide with~$c_1$. Source~(iv) is the only one intrinsically
independent of $d_0$, and the Class~4 analysis below shows only the
identity $f(\sqrt{\pi}) = \sqrt{\pi}$ produces a primitive.

\emph{Class~1: Values of $p$ and its decomposition at $d_0$.} The zero
gives three scalars: $p(d_0) = 0$, the $d$-dependent part
$q(d_0) = \frac{1}{2}\ln\pi$, and the constant part
$-\frac{1}{2}\ln\pi$. (i)~$p(d_0) = 0$ fails~(a): $p(d) = 0$ has its
solution at $d_0$ itself. (ii)~$q(d_0) = \frac{1}{2}\ln\pi$ satisfies~(a)
since $\frac{1}{2}\ln\pi > 0$, and satisfies~(b) since
$\frac{1}{2}\ln\pi$ is the log-space image of $\sqrt{\pi}$, the
recurrence's unique primitive. This is $c_1$. (iii)~$-\frac{1}{2}\ln\pi < 0$
fails~(a).

\emph{Class~2: Higher derivatives $p^{(n)}(d_0)$.} All fail~(b): they are
not recurrence primitives.

\emph{Class~3: The dimension $d_0$ itself.} $p(d) = d_0$ equates a
dimensionless rate to a dimension count---a type mismatch.

\emph{Class~4: Functions of $c_1$ (equivalently $\sqrt{\pi}$).} Consider
$f(c_1)$ for any $f\neq\mathrm{id}$.
%
\emph{Multiples:} $n c_1 = \tfrac{1}{2}\ln(\pi^n)$ is the constant of
the recurrence for $V_d^{\,n}$, a derived sequence. A primitive of a
derived sequence is not a primitive of the original first-order recurrence.
%
\emph{Powers:} $c_1^{\,k}$ for $k\geq 2$ does not appear as the
dimension-independent constant in any form of the recurrence; it is
algebraically independent of both $\sqrt{\pi}$ and $\frac{1}{2}\ln\pi$.
%
\emph{General:} any analytic $f(c_1)\neq c_1$ either reduces to a
multiple $n c_1$ (hence derived) or produces a scalar that is not a
recurrence constant in any canonical form of the first-order recurrence.

Across all four classes, $c_1 = \frac{1}{2}\ln\pi$ is unique.
\end{proof}

\begin{remark}[Alternative candidates addressed]\label{rem:c1-alternatives}
Four natural candidates not written as members of Classes~1--3 above are
scalars built from the derived cascade geometry at $d_0$:
$R(d_0)$, the area ratio $\Omega_{d_0-1}/\Omega_{d_0+1}$, the Stirling
expression $2\pi e^{2\sqrt{\pi}}$ (appearing in Theorem~\ref{thm:invariant}
below), and the downstream primitive $2\sqrt{\pi}$
(Corollary~2.3 of Paper~IVb). Each fails the primitivity criterion~(b),
by one of two routes: the first two depend nontrivially on $d_0$ and
fall outside Class~4 (excluded by the dimension-independence argument
of Theorem~\ref{thm:c1unique}); the latter two lie in Class~4 but are
multiples of the primitive, not the primitive itself:
\begin{itemize}
\item $R(d_0) = \Gamma((d_0{+}1)/2)/\Gamma((d_0{+}2)/2)$ is a
\emph{derived} cascade quantity (Corollary~\ref{cor:primary}), not a
recurrence primitive; it depends on $d_0$ and on $\sqrt{\pi}$ together,
not on a single dimension-independent constant.
\item $\Omega_{d_0-1}/\Omega_{d_0+1} = 1/(\pi\,R(d_0-1)\,R(d_0))$ is a
rational multiple of $1/\pi$ depending on $d_0$; it is a derived ratio of
sphere areas, not a recurrence constant.
\item $2\pi\,e^{2\sqrt{\pi}}$ is an analytic combination of
$\sqrt{\pi}$ and $\pi$ and fails~(b) for the same reason as Class~4
\emph{powers}: not the dimension-independent constant of any canonical
recurrence form.
\item $2\sqrt{\pi} = N(0)\cdot\Gamma(\tfrac{1}{2})$
(Part~IVb, Corollary~\xref{part4b}{cor:2sqrtpi-primitive}) is a \emph{product} of the primitive
$\sqrt{\pi}$ with a derived cascade quantity ($N(0) = 2$, the $d{=}0$
slicing lapse); as a multiple of the primitive it is a derived scalar, not
itself the primitive.
\end{itemize}
No candidate in the SP-5 list is a primitive of the first-order
recurrence: the first two are excluded by $d_0$-dependence; the latter
two are Class~4 but fail to be \emph{the} dimension-independent
multiplicative constant of the recurrence.
\end{remark}

\subsection{Two canonical values of one constant}

The recurrence contains one irreducible constant
(Theorem~\ref{thm:irred}). That constant has two canonical
representations, because the recurrence has two canonical forms:

\begin{center}
\begin{tabular}{lll}
\toprule
Form & Expression & Constant \\
\midrule
Multiplicative & $V_{d+1}/V_d = \sqrt{\pi}\cdot R(d+1)$ &
$c_2 = \sqrt{\pi}$ \\
Logarithmic & $\Delta\ln V_d = \frac{1}{2}\ln\pi + \ln R(d+1)$ &
$c_1 = \frac{1}{2}\ln\pi$ \\
\bottomrule
\end{tabular}
\end{center}

These are the same constant in two algebraic spaces, related by the
exponential map: $e^{c_1} = e^{\frac{1}{2}\ln\pi} = \sqrt{\pi} = c_2$.

\begin{definition}[Canonical form]\label{def:canonical}
A canonical form of a first-order recurrence is an algebraically equivalent
expression that isolates the dimension-independent content as a single
constant with coefficient~1. For a multiplicative recurrence
$a_{n+1} = K\cdot f(n)\cdot a_n$, the two canonical forms are the
recurrence itself (constant $K$) and its logarithm (constant $\ln K$).
Any other algebraically equivalent expression either has a constant that
is a power, root, or inverse of one of these (hence derived, not
primitive) or fails to isolate the constant with coefficient~1.
\end{definition}

\begin{remark}[Why coefficient~1]\label{rem:coeff1}
The condition ``coefficient~1'' is not a convention; it is forced by the
recurrence's order. The slicing recurrence is first-order: each step
applies $K$ exactly once. A form containing $K^n$ ($n\neq 1$) describes
one of two things: either the recurrence for a derived sequence
($V_d^{\,n}$ satisfies $V_{d+1}^{\,n}/V_d^{\,n} = K^n\cdot R(d+1)^n$), or
the $n$-step composition $V_{d+n}/V_d$, which is a higher-order object.
In neither case is $K^n$ a primitive of the original first-order
recurrence. Excluding non-unit exponents is therefore not a filter designed
to produce a desired output---it is the statement that the paper studies
the recurrence as given, not a power or iterate of it.
\end{remark}

\begin{theorem}[Two canonical values]\label{thm:twovalues}
The recurrence's dimension-independent content is exactly the pair
$(c_1, c_2) = (\tfrac{1}{2}\ln\pi,\;\sqrt{\pi})$: one constant in each
canonical form. No third canonical value exists.
\end{theorem}

\begin{proof}
By Definition~\ref{def:canonical}, a first-order multiplicative recurrence
has exactly two canonical forms. The multiplicative form isolates
$\sqrt{\pi}$; the logarithmic form isolates $\frac{1}{2}\ln\pi$. Any
other rewriting---squaring ($\pi$), inverting ($1/\sqrt{\pi}$), taking
roots ($\pi^{1/4}$)---produces a derived constant, not a primitive one:
the derived constant is a function of $\sqrt{\pi}$ and cannot appear with
coefficient~1 in any first-order form of the recurrence. For the slicing
recurrence: $K = \sqrt{\pi}$, $\ln K = \frac{1}{2}\ln\pi$. These exhaust
the dimension-independent content.
\end{proof}

\begin{remark}
The first threshold comes from applying $p(d) = c$ to the logarithmic
constant $c_1 = \frac{1}{2}\ln\pi$. The second comes from applying the
same construction to the multiplicative constant $c_2 = \sqrt{\pi}$.
Using only one would ignore half the recurrence's constant content. The
bridge between them is not an arbitrary operation: it is the exponential
map that defines the relationship between a multiplicative recurrence and
its logarithm.
\end{remark}

\subsection{Orbit termination}

\begin{theorem}[Orbit termination]\label{thm:terminate}
The orbit of $\{0\}$ under the threshold construction is
$\{0, c_1, c_2\}$. No further element is reachable.
\end{theorem}

\begin{proof}
The orbit is generated by two operations: (A)~reading $c_1$ from the
natural zero, and (B)~passing to the other canonical form via $\exp$.
A first-order multiplicative recurrence has exactly two canonical forms
(Theorem~\ref{thm:twovalues}), connected by exactly one bridge map
($\exp/\ln$). The orbit visits the zero (seed), the logarithmic constant
($c_1$), and the multiplicative constant ($c_2$). No fourth station
exists because there is no third canonical form.

Explicitly: (i)~$\exp(c_2) = e^{\sqrt{\pi}}\approx 5.83$ is not a
recurrence primitive by Theorem~\ref{thm:irred}---the recurrence's
only nontrivial constant is $\sqrt{\pi}$, and $e^{\sqrt{\pi}}$ does not
appear. (ii)~$\ln(c_1) = \ln(\frac{1}{2}\ln\pi) \approx -0.558$ is
negative and fails to generate a threshold beyond $d_0$.
(iii)~Operation~A cannot be reapplied: $p$ has exactly one zero (strict
monotonicity), so there is no second zero at which to read off a new
value.

The orbit terminates at three elements because the recurrence has two
canonical values and one structural zero. This is a property of the
recurrence's order, not a coincidence about specific numbers.
\end{proof}

\subsection{Two thresholds and no more}

\begin{theorem}[Two thresholds]\label{thm:thresholds}
$p(d) = c_1$ has a unique continuous solution at
$d_1^\star \approx 19.7308$, lying between integers $19$ and $20$.
$p(d) = c_2$ has a unique continuous solution at
$d_2^\star \approx 217.6267$, lying between integers $217$ and $218$.
The canonical integer labels $d_1 = 19$ and $d_2 = 217$ are fixed in
Theorem~\ref{thm:variational} below by the variational (argmax)
characterisation of the cascade invariant.
\end{theorem}

\begin{proof}
Strict monotonicity of $p$ gives uniqueness of the continuous solutions.
Discrete verification:
$p(19) = 0.55351 < c_1 = 0.57236 < p(20) = 0.57914$.
$p(217) = 1.77101 < c_2 = 1.77245 < p(218) = 1.77331$.
Both continuous solutions lie strictly between the bracketing integers.
\end{proof}

\begin{corollary}[No third threshold]\label{cor:no-third-threshold}
The orbit terminates at $(c_1, c_2)$. Each positive canonical value
determines exactly one threshold dimension. The cascade has exactly two
thresholds.
\end{corollary}

\begin{proof}
Theorem~\ref{thm:terminate} establishes that the orbit of $\{0\}$ under
the threshold construction is exactly $\{0, c_1, c_2\}$, so no further
canonical value exists. By Theorem~\ref{thm:thresholds}, each positive
canonical value $c\in\{c_1, c_2\}$ gives a unique threshold dimension
via strict monotonicity of $p$. Two canonical values, two thresholds.
\end{proof}

% ============================================================
\section{Four Distinguished Dimensions}\label{sec:four-distinguished-dimensions}
% ============================================================

\begin{theorem}[Tower completeness]\label{thm:tower}
The Gamma function produces exactly four distinguished dimensions in the
cascade. No fifth exists.
\end{theorem}

\begin{proof}
\emph{Cascade-internal functions of dimension are analytic in
$(\sqrt{\pi}, d)$.} By Theorem~\ref{thm:irred} and
Corollary~\ref{cor:primary}, every cascade-internal function of
dimension---$\Omega_d$, $V_d$, $R(d)$, $p(d)$, $R_{\mathrm{eff}}(d)$
(Theorem~\ref{thm:reff}), and analytic combinations thereof---is an
analytic function $G(d) = F(\sqrt{\pi}, d)$ of the recurrence's
primitive $\sqrt{\pi}$ and the dimension $d$ (extended to
$\mathbb{R}_{\geq 0}$ via the Gamma function). This is the
function-valued analogue of Theorem~\ref{thm:c1unique}'s closure,
established by the same chain of theorems: sphere area is the unique
independent cascade quantity (Corollary~\ref{cor:primary}), the
recurrence has a unique primitive $\sqrt{\pi}$
(Theorem~\ref{thm:irred}), and every other cascade quantity is an
analytic function of $\sqrt{\pi}$ and $d$.

\emph{Integer-selecting operators on cascade-internal functions.} A
cascade-internal procedure that selects an integer dimension consists
of choosing such a function $G$ and applying an integer-selecting
operator to it. The integer-selecting operators on an analytic
real-valued function of one variable that yield a single distinguished
integer are:
\begin{itemize}
\item[(a)] \emph{Extremization.} $d^* = \operatorname{argmax} G$ (or
$\operatorname{argmin}$). Yields an interior integer only when $G$ is
unimodal; for strictly monotonic $G$, the extremum is at the boundary
and selects no interior integer.
\item[(b)] \emph{Sign change (zero).} $d^* = G^{-1}(0)$. Defined when
$G$ takes both signs; unique when $G$ is strictly monotonic.
\item[(c)] \emph{Level crossing with a primitive reference.}
$d^* = G^{-1}(c)$ for $c$ a cascade-primitive constant
(Theorem~\ref{thm:c1unique}, Class~4). The dimension-independence of
$c$ is essential: a $d$-dependent reference $c(d)$ reduces the equation
$G(d) = c(d)$ to the zero of $G - c$, hence class~(b).
\item[(d)] \emph{Higher-order critical points.} $d^* = (G^{(n)})^{-1}(0)$
for some derivative order $n \geq 1$ (zeros of derivatives, inflection
points, etc.).
\end{itemize}

\emph{No fifth operator class.} Any further integer-selecting
construction reduces to one of (a)--(d):
\begin{itemize}
\item \emph{Composition.} $\operatorname{argmax} G^{(n)}$ is the zero
of $G^{(n+1)}$, hence class~(d).
\item \emph{Integer offsets} from a class-(c) selection ($d_c + k$ for
fixed integer $k$) are reachable from $d_c$ by the recurrence's own
stepping operation (Theorem~\ref{thm:c1unique}'s integer-arithmetic
closure), not independent selections.
\item \emph{Intersections} of two cascade-internal functions, $G_1(d) =
G_2(d)$, reduce to the zero of $G_1 - G_2$, hence class~(b).
\item \emph{Asymptotic invariants} (Lyapunov exponents, growth rates,
limit-set indices) produce real numbers rather than specific integer
dimensions; converting them to integer outputs (via thresholding or
rounding) reduces to a level-crossing equation, hence class~(c) or~(b).
\end{itemize}
The classes~(a)--(d) therefore exhaust cascade-internal
dimension-selecting mechanisms. Theorem~\ref{thm:c1unique}'s
primitivity criterion provides the additional constraint on class~(c):
no non-primitive reference value introduces a new distinguished
dimension.

\emph{Classes~(a) and (b) jointly yield two dimensions.}
Two equilibria: $V_d$ has a unique maximum at $d_V = 5$ (by unimodality;
class~(a) applied to $V_d$), and $p(d)$ has a unique zero at $d_0 = 7$
(by strict monotonicity; class~(b) applied to $p$). The only other
cascade sequence with a local extremum is $\Omega_d$, whose argmax
$d_{\max} = 6$ is \emph{derived}: by Theorem~\ref{thm:dual},
$d_{\max} = d_0 - 1$, so $d = 6$ is a derived consequence of $d_0 = 7$
via boundary dominance. Including it would double-count the decay-rate
zero. The volume and sphere-area arg\-max thus reduce to a single
independent class-(a) selection, $d_V = 5$.

\emph{Class~(c) yields two thresholds.}
The complete cascade-primitive constant content generates exactly
$(c_1, c_2)$---$c_1 = \tfrac{1}{2}\ln\pi$ is unique by
Theorem~\ref{thm:c1unique}, $c_2 = \sqrt{\pi}$ is unique as the
multiplicative primitive of the recurrence
(Theorem~\ref{thm:irred})---and the orbit terminates
(Theorem~\ref{thm:terminate}), giving $d_1 = 19$ and $d_2 = 217$. No
further cascade-primitive reference values exist: any additional
candidate reference would either be a multiple of $c_1$ or $c_2$
(excluded by Theorem~\ref{thm:c1unique}'s Class~4) or a non-cascade
import (excluded by the ``cascade-internal'' qualifier; see
Remark~\ref{rem:tower-alternatives} for specific cases).

\emph{Class~(d) contributes no new dimensions.}
Higher derivatives $p^{(n)}$ fail Theorem~\ref{thm:c1unique}'s
primitivity criterion (Class~2 in that proof): they are not recurrence
primitives, and their zeros are not canonical cascade dimensions.
Inflection points of $\Omega_d$ reduce to zeros of $p'$, which are zeros
of a non-primitive quantity---same exclusion.

No other cascade quantity selects a dimension.
\end{proof}

\begin{remark}[Alternative candidates addressed]\label{rem:tower-alternatives}
Three alternative mechanisms that could, in principle, select a fifth
distinguished dimension are not cascade-internal and are therefore
excluded by Theorem~\ref{thm:tower}:
\begin{itemize}
\item \emph{Crossings with the self-dual radius $R(d) = 1/\sqrt{2}$.}
The equation has no integer solution: $R(11) \approx 0.7520$,
$R(12) \approx 0.7087$, $R(13) \approx 0.7201$, so the continuous
crossing lies between $d=12$ and $d=13$ and no integer satisfies the
equation exactly. Moreover $1/\sqrt{2}$ is not a cascade primitive: it
arises as a self-dual radius in Kaluza--Klein compactification, a
procedure the series explicitly refuses (Paper~I~\S3.2). Part~IVa's
Theorem~2.1 notes a 0.225\% near-miss at $d=12$ (documented as a
separate soft spot, SP-27, and flagged there as a Check-7 tension); the
Tower Completeness theorem here is not disturbed by the near-miss
because no exact cascade-internal mechanism produces $d=12$ as a
distinguished dimension.
\item \emph{Zeros of higher derivatives $p^{(n)}$ for $n\geq 1$.}
$p'(d) = \tfrac{1}{4}\psi^{(1)}((d+1)/2) > 0$ everywhere (no zeros),
$p''(d) = \tfrac{1}{8}\psi^{(2)}((d+1)/2) < 0$ everywhere (no zeros). So
higher derivatives of $p$ have no zeros at any dimension---the
monotonicity class is already captured by $p$ itself.
\item \emph{Inflection points of $\Omega_d$.} Reduce to zeros of the
first derivative of the log-area, which is $-p(d)$. The only zero of
$p$ is $d_0 = 7$; the corresponding inflection in $\Omega_d$ is already
encoded in $d_0$ and is not an independent selection.
\end{itemize}
No fifth cascade-internal mechanism exists.
\end{remark}

\begin{center}
\begin{tabular}{ccrl}
\toprule
Index & $d$ & $\Omega_d$ & Role \\
\midrule
0 & 5 & $\pi^3$ & Volume maximum (argmax $V_d$) \\
1 & 7 & $\pi^4/3$ & Decay-rate zero (natural zero of $p$) \\
2 & 19 & 0.516 & First threshold ($c_1 = \tfrac{1}{2}\ln\pi$) \\
3 & 217 & $2.335\times 10^{-120}$ &
Second threshold ($c_2 = \sqrt{\pi}$) \\
\bottomrule
\end{tabular}
\end{center}

\subsection{The regime partition}\label{sec:regimes}

The three privileged values of $p(d)$---zero, $c_1 = \tfrac{1}{2}\ln\pi$,
and $c_2 = \sqrt{\pi}$---partition the cascade into exactly four
qualitatively distinct regimes. Since $p(d)$ is strictly monotone
increasing (Theorem~\ref{thm:dual}) and crosses each privileged value
exactly once, the partition is well-defined:

\begin{center}
\renewcommand{\arraystretch}{1.2}
\begin{tabular}{llll}
\toprule
Regime & Condition on $p(d)$ & Integer range & Character \\
\midrule
Growth & $p(d) < 0$ & $d \leq 6$ & $\Omega_d$ increasing \\
Decay & $0 \leq p(d) < c_1$ & $7 \leq d \leq 19$ &
  Subcritical: per-step ratio $< \sqrt{\pi}$ \\
Exponential decay & $c_1 \leq p(d) < c_2$ & $20 \leq d \leq 217$ &
  Supercritical: per-step ratio $> \sqrt{\pi}$ \\
Oblivion & $p(d) \geq c_2$ & $d \geq 218$ &
  $\Omega_d < 10^{-120}$; geometrically null \\
\bottomrule
\end{tabular}
\end{center}

\noindent The four distinguished dimensions of Theorem~\ref{thm:tower}
separate into two structurally distinct categories under this partition:

\begin{itemize}
\item \textbf{Regime boundaries} (three of the four). The integers
$d_0 = 7$, $d_1 = 19$, and $d_2 = 217$ are selected by the three
continuous crossings of $p(d)$. Each continuous crossing lies strictly
between integers: $d_0^\star \approx 6.257$, $d_1^\star \approx 19.731$,
$d_2^\star \approx 217.627$. The integer labels assigned here are
fixed by the variational characterisation of the cascade invariant
(Theorem~\ref{thm:variational} below), not by an ad hoc rounding
convention.
\item \textbf{Interior landmark} (one of the four). $d_V = 5$ is the
discrete argmax of $V_d$, sitting inside the Growth regime. It is one
index below the discrete argmax of $\Omega_d$ at $d = 6$, and both
landmarks are interior to Growth---neither is a regime boundary. The
observer at $d = 4$ resides on the $S^3$ boundary of $B^5$ (the volume
maximum ball), also inside Growth.
\end{itemize}

The regime partition clarifies the architecture that the
four-distinguished-dimension list of Theorem~\ref{thm:tower} compresses:
each regime is bounded below and above by one of the three crossings
(except Growth, which extends to $d=0$, and Oblivion, which extends to
$d=\infty$), and the four distinguished dimensions record one interior
landmark plus the three boundaries.

% ============================================================
\section{The Cascade Invariant}\label{sec:the-cascade-invariant}
% ============================================================

Theorem~\ref{thm:tower} establishes exactly four distinguished dimensions.
Two are thresholds: $d_1 = 19$ and $d_2 = 217$, selected by the canonical
values of the recurrence constant. Two are equilibria: $d_V = 5$ (volume
maximum) and $d_0 = 7$ (decay-rate zero). These four exhaust the Gamma
function's distinguished structure.

The cascade invariant is the unique scalar built from these four.

\subsection{Content and scale from structural stability}

The four distinguished dimensions are selected by two different
mechanisms. This subsection proves the mechanisms determine their algebraic
role in the invariant.

\begin{definition}[Generalised recurrence]\label{def:general}
Let $K > 0$ be any constant. The generalised slicing recurrence is
$V_{d+1} = V_d\cdot K\cdot R(d+1)$, with decay rate
$p_K(d) = -\ln K + \tfrac{1}{2}\psi((d+1)/2)$. The physical cascade has
$K = \sqrt{\pi}$, but the structural analysis applies for all $K$.
\end{definition}

\begin{theorem}[Structural stability of the equilibria]\label{thm:stable}
For any $K > 0$, the generalised recurrence has:
\begin{enumerate}
\item[(a)] a unique volume maximum $d_V(K)$, where $K\cdot R(d) = 1$;
\item[(b)] a unique decay-rate zero $d_0(K)$, where $p_K(d) = 0$.
\end{enumerate}
Their existence is guaranteed by monotonicity of $R(d)$ and $\psi(d)$
respectively. Their locations shift continuously with $K$, but they
exist for every $K > 0$.
\end{theorem}

\begin{proof}
(a)~$R(d) = \Gamma((d+1)/2)/\Gamma((d+2)/2)$ is positive, continuous, and
strictly decreasing with $R(d)\to 0$ as $d\to\infty$. For any $K > 0$,
the equation $K\cdot R(d) = 1$ has a unique solution by the intermediate
value theorem.

(b)~$\psi((d+1)/2)$ is strictly increasing and unbounded. Therefore
$p_K(d) = -\ln K + \frac{1}{2}\psi((d+1)/2)$ crosses zero exactly once
for any $K > 0$.
\end{proof}

\begin{theorem}[Structural sensitivity of the thresholds]\label{thm:sensitive}
The generalised recurrence's two canonical values are $c_1(K) = \ln K$
and $c_2(K) = K$. The threshold dimensions $d_1(K)$ and $d_2(K)$, defined
by $p_K(d_1) = c_1(K)$ and $p_K(d_2) = c_2(K)$, depend on the specific
value of $K$:

\begin{center}
\begin{tabular}{cccc}
\toprule
$K$ & $d_0$ & $d_1$ & $d_2$ \\
\midrule
$1.50$ & 5 & 11 & 91 \\
$\sqrt{\pi}$ & 7 & 19 & 217 \\
$2.00$ & 8 & 32 & 437 \\
$2.50$ & 13 & 79 & 1856 \\
\bottomrule
\end{tabular}
\end{center}

The threshold dimensions, and the sphere areas at those dimensions, are
determined by the specific numerical value of $K$. Different $K$, different
thresholds, different content.
\end{theorem}

\begin{theorem}[Content--scale separation]\label{thm:cs}
The four distinguished dimensions separate into two classes by structural
stability:

\emph{Scale dimensions} ($d_V$, $d_0$): selected by the dynamical
conditions $K\cdot R(d) = 1$ and $p_K(d) = 0$. These conditions ask where
the recurrence constant \emph{balances} the dimension-dependent factor.
They exist for all $K > 0$ (Theorem~\ref{thm:stable}). They are
structurally stable: properties of the recurrence's dynamics, not its
constant's value.

\emph{Content dimensions} ($d_1$, $d_2$): selected by the conditions
$p_K(d) = c_1(K)$ and $p_K(d) = c_2(K)$. These conditions ask where the
decay rate \emph{equals} the recurrence constant. They depend on the
specific value of $K$ (Theorem~\ref{thm:sensitive}). They are structurally
sensitive: properties of the recurrence's constant, not its dynamics.

In the invariant:
\begin{enumerate}
\item[(i)] Scale dimensions enter as a \emph{ratio}. Two equilibria define
two reference scales; their relationship is a comparison:
$\Omega_{d_V}/\Omega_{d_0}$.

\item[(ii)] Content dimensions enter as a \emph{product}. Two level
crossings define two independent outputs of the recurrence at
characteristic rates; independent outputs of a multiplicative process
combine multiplicatively: $\Omega_{d_1}\times\Omega_{d_2}$.
\end{enumerate}
\end{theorem}

\begin{proof}
(i)~The scale dimensions are selected by equilibrium conditions: at $d_V$,
the volume recurrence ratio $K\cdot R(d)$ balances to unity; at $d_0$,
the decay rate $p(d)$ crosses zero. Both mark transitions in the
recurrence's dynamics---growth-to-decay for volumes at $d_V$,
growth-to-decay for sphere areas near $d_0$. The sphere areas at these
transitions serve as reference scales: they measure the recurrence near
its turning points, not at characteristic values of the constant.
Two equilibria give two references. Their relationship is a ratio:
$\Omega_{d_V}/\Omega_{d_0}$.

(ii)~At a canonical level crossing, the decay rate matches a value derived
from the recurrence constant. The sphere area at this crossing is a
specific output: how much sphere area remains after the decay rate has
reached the characteristic value $c_1$ or $c_2$. Two level crossings at
independent canonical values give two independent outputs. By the
recurrence's multiplicative structure
(Corollary~\ref{cor:primary}), independent outputs combine as a product:
$\Omega_{d_1}\times\Omega_{d_2}$.

Scale dimensions cannot be content: an equilibrium condition marks a
dynamical transition, which exists for any $K$. Including it as a
content factor would make the invariant depend on a structural feature
that carries no information about the specific constant $\sqrt{\pi}$.
Content dimensions cannot be scale: a canonical crossing depends on the
specific value of $K$; using it as a reference would make the scale
depend on the content.
\end{proof}

\begin{remark}[The separation is not a definition]
Theorem~\ref{thm:cs} is a consequence of structural stability, not a
classification imposed by hand. The equilibria are scale dimensions because
they survive changes to $K$; the thresholds are content dimensions because
they don't. If the Gamma function had a different structure---say, if
$\psi$ were not monotone and $p$ had multiple zeros---the separation might
not hold. It holds here because the slicing recurrence is first-order with
a strictly monotone dimension-dependent factor. This is a theorem about
the recurrence's order and monotonicity.
\end{remark}

\begin{remark}[Why scale enters as a ratio, not a product: frame conversion]\label{rem:sp7-status}
The assignment ``scale = ratio'' in part~(i) of Theorem~\ref{thm:cs} is
not a stylistic choice between two equally valid scalars; it is the
unique frame-conversion factor forced by the first-order bilinear
structure of the recurrence. The argument has three steps, each
provable from the theorems already in hand:
\begin{enumerate}
\item[(1)] \emph{Bilinear decomposition
(Lemma~\ref{lem:bilinear}).} Because the recurrence is first-order
multiplicative, every sphere area decomposes as
$\Omega_d = \Omega_{d_{\mathrm{ref}}}\cdot\mathcal{R}_d$ with the
reference factor $\Omega_{d_{\mathrm{ref}}}$ appearing at exponent
exactly~$1$. A product of $n$ sphere areas therefore contains
$\Omega_{d_{\mathrm{ref}}}^{\,n}$.
\item[(2)] \emph{Intrinsic reference.} The natural intrinsic reference
is $d_0$, the decay-rate zero: $p(d_0) = 0$ is the unique place where
the recurrence's constant $c_1 = \tfrac{1}{2}\ln\pi$ balances its
dimension-dependent factor, so $\Omega_{d_0}$ is the cascade's own
normalisation scale. The content product
$\Omega_{d_1}\cdot\Omega_{d_2}$ implicitly carries
$\Omega_{d_0}^{\,2}$ under this decomposition.
\item[(3)] \emph{Observer's reference and frame conversion.} The
observer lives at the $S^3$ boundary of $B^5$, so the observer-frame
reference is $\Omega_{d_V}$ rather than $\Omega_{d_0}$. Replacing
$\Omega_{d_0}^{\,2}$ with $\Omega_{d_V}^{\,2}$ in the content product
requires multiplication by the frame-conversion factor
$(\Omega_{d_V}/\Omega_{d_0})^{2}$. This is the scale term in the
invariant: it is a \emph{ratio} because it is a conversion from one
reference to another, and it has exponent~$2$ because the content
product carries two implicit reference factors
(Lemma~\ref{lem:bilinear}).
\end{enumerate}
A product of the scale dimensions ($\Omega_{d_V}\cdot\Omega_{d_0}$)
cannot substitute for this ratio: it would introduce an additional
factor of $\Omega_{d_V}$ unrelated to the observer's frame conversion,
violating the exponent-counting rule of
Lemma~\ref{lem:bilinear}. Conversely, a ratio of the content
dimensions ($\Omega_{d_1}/\Omega_{d_2}$) would violate
Theorem~\ref{thm:content}: in the $d_2\to\infty$ limit, the ratio
diverges while the physical content must vanish. Each operation is
fixed by a different forcing argument---product for content by the
boundary-limit behaviour of Theorem~\ref{thm:content}, ratio for
scale by the frame-conversion argument above---and neither is a
definitional choice.
\end{remark}

\subsection{The complete invariant}

\begin{theorem}[Content, not scale]\label{thm:content}
The natural combination of the threshold sphere areas is their product,
not their quotient.
\end{theorem}

\begin{proof}
The cascade originates at $d = \infty$ where
$V_\infty = \Omega_\infty = 0$: geometric nothing. As
$d_2\to\infty$: the product $\Omega_{d_1}\times\Omega_{d_2}\to 0$
(content approaches nothing---correct). The quotient
$\Omega_{d_1}/\Omega_{d_2}\to\infty$ (scale diverges at the origin---wrong).

The recurrence is first-order multiplicative
(Corollary~\ref{cor:primary}). Independent contributions at independent
levels combine as a product.
\end{proof}

\begin{theorem}[Two scales]\label{thm:twoscales}
The ratio $\Omega_{d_V}/\Omega_{d_0} = \Omega_5/\Omega_7 = 3/\pi$ is the
unique dimensionless scalar relating the cascade's two equilibria.
\end{theorem}

\begin{proof}
By Corollary~\ref{cor:primary}, sphere area is the unique independent
cascade quantity at each level. The ratio of sphere areas at the two
equilibria is therefore the unique dimensionless relationship between them.
Evaluation:
$\Omega_5 = \pi^3$,\;
$\Omega_7 = \pi^4/3$,\;
ratio $= 3/\pi$.
\end{proof}

\begin{lemma}[Scale ordering]\label{lem:ordering}
$\Omega_{d_V} < \Omega_{d_0}$.
\end{lemma}

\begin{proof}
Direct computation: $\Omega_5 = \pi^3\approx 31.0$,\;
$\Omega_7 = \pi^4/3\approx 32.5$.
More generally, the volume peak $d_V$ lies below the decay-rate zero
$d_0$ because $V_d = \Omega_{d-1}/d$ shifts the volume equilibrium to a
lower index. Both lie within the near-peak region of $\Omega_d$
(whose discrete maximum is at $d = 6$; see Theorem~\ref{thm:dual}).
The ordering $\Omega_{d_V} < \Omega_{d_0}$ holds because $d_V = 5$
is further below the discrete peak than $d_0 = 7$ is above it, and
the log-convexity of $1/\Omega_d$ ensures steeper growth than decay
near the peak.
\end{proof}

\begin{corollary}\label{cor:direction}
$\Omega_{d_V}/\Omega_{d_0} < 1$, and $(\Omega_{d_V}/\Omega_{d_0})^n < 1$
for any $n\geq 1$. In the physical cascade: $\Omega_5/\Omega_7 = 3/\pi
< 1$.
\end{corollary}

\begin{lemma}[First-order reference decomposition]\label{lem:bilinear}
In the first-order multiplicative recurrence, any sphere area $\Omega_d$
decomposes as
$\Omega_d = \Omega_{d_0}\cdot\mathcal{R}_d$,
where $\mathcal{R}_d := \prod_{k=d_0}^{d-1}(\Omega_{k+1}/\Omega_k)$ is
the intrinsic ratio (product of step factors, independent of any reference
choice) and $\Omega_{d_0}$ appears at exponent exactly~1. The exponent is
forced: the recurrence outputs one $\Omega$ per step.

Consequently, a product of $n$ sphere areas contains
$\Omega_{d_0}^{\,n}$:
\[
\prod_{i=1}^n \Omega_{d_i}
= \Omega_{d_0}^{\,n}\prod_{i=1}^n \mathcal{R}_{d_i}.
\]
A scale correction replacing $\Omega_{d_0}$ with $\Omega_{d_V}$ therefore
appears at exponent~$n$.
\end{lemma}

\begin{proof}
Each $\Omega_{d_i} = \Omega_{d_0}\cdot\mathcal{R}_{d_i}$ by telescoping
the recurrence. The product collects $n$ factors of $\Omega_{d_0}$.
\end{proof}

\begin{theorem}[Cascade invariant]\label{thm:invariant}
The unique dimensionless scalar built from all four distinguished sphere
areas, under the constraints of the recurrence, is
\[
I = \frac{\Omega_5^2}{\Omega_7^2}\;\Omega_{19}\;\Omega_{217}
= \frac{9}{\pi^2}\;\Omega_{19}\;\Omega_{217}.
\]
\end{theorem}

\begin{proof}
Six constraints determine the invariant.

\emph{(1)~Sphere areas.} By Corollary~\ref{cor:primary}, $\Omega_d$ is the
unique independent cascade quantity. All factors are sphere areas.

\emph{(2)~All four distinguished dimensions.}
Theorem~\ref{thm:tower} proves exactly four distinguished dimensions.
No cascade mechanism selects a proper subset of these four: the equilibria
are selected by dynamical transitions ($K\cdot R(d) = 1$ and $p = 0$),
the thresholds by canonical values ($p = c_1$ and $p = c_2$), and these
are independent selection mechanisms with no hierarchy between them.
An invariant omitting any distinguished dimension discards information the
recurrence provides; an invariant including all four uses the complete
output. Completeness is not imposed---it is the absence of an arbitrary
omission: $\Omega_5$, $\Omega_7$, $\Omega_{19}$, $\Omega_{217}$.

\emph{(3)~Threshold exponents.} The recurrence is first-order
multiplicative, outputting one $\Omega$ per step. Each threshold
contributes one sphere area at exponent~1. Two thresholds: exponents
$c = d = 1$.

\emph{(4)~Product for content.} The content dimensions (thresholds) combine
as a product (Theorem~\ref{thm:content}):
$\Omega_{19}\times\Omega_{217}$.

\emph{(5)~Ratio for scale.} The scale dimensions (equilibria) combine as a
ratio (Theorem~\ref{thm:cs}): $\Omega_5/\Omega_7$.

\emph{(6)~Scale exponent and direction.} The content product has exactly
two factors (one per threshold, exponent~1 each by constraint~3).
The exponent on the scale ratio equals the number of content factors:
$a = 2$.  This follows from Lemma~\ref{lem:bilinear} applied to any
reference dimension $d_{\mathrm{ref}}$: each sphere area $\Omega_d$
decomposes as $\Omega_{d_{\mathrm{ref}}}\cdot\mathcal{R}_d$, with
$\Omega_{d_{\mathrm{ref}}}$ at exponent exactly~1 (the recurrence is
first-order). A product of $n$ sphere areas therefore contains
$\Omega_{d_{\mathrm{ref}}}^{\,n}$. The exponent $n$ counts the
factors---it does not depend on the reference choice. Since orbit
termination (Theorem~\ref{thm:terminate}) produces exactly two thresholds,
$n = 2$.  The scale correction
$(\Omega_{d_V}/\Omega_{d_0})^2$ then converts from any reference to the
observer's equilibrium. The direction
$\Omega_{d_V}/\Omega_{d_0} < 1$ is forced by the scale ordering
(Lemma~\ref{lem:ordering}).

The invariant is $(\Omega_5/\Omega_7)^2\times\Omega_{19}\times\Omega_{217}
= \Omega_5^2\,\Omega_{19}\,\Omega_{217}/\Omega_7^2$. The exponents are
$a = 2$, $b = -2$, $c = 1$, $d = 1$, uniquely determined.
\end{proof}

\begin{remark}[Burden of uniqueness]\label{rem:burden}
The invariant is unique given six constraints. Each constraint traces to a
theorem: sphere-area primacy (Corollary~\ref{cor:primary}), tower
completeness (Theorem~\ref{thm:tower}), first-order exponents, the product
rule (Theorem~\ref{thm:content}), the ratio rule
(Theorem~\ref{thm:cs}), and the bilinear decomposition
(Lemma~\ref{lem:bilinear}). No constraint is an axiom: each is a proved
consequence of the recurrence's order, monotonicity, and
irreducible constant. The conjunction of six proved statements is itself
proved.
\end{remark}

\begin{theorem}[Evaluation]
$I = 1.0989\times 10^{-120}$.
\end{theorem}

\begin{proof}
$\Omega_5 = \pi^3$.\;
$\Omega_7 = \pi^4/3$.\;
$\Omega_{19} = 0.51614$.\;
$\Omega_{217} = 2.33490\times 10^{-120}$.\;
$(\Omega_5/\Omega_7)^2 = 9/\pi^2 = 0.91189$.\;
Product: $0.91189\times 0.51614\times 2.33490\times 10^{-120}
= 1.09895\times 10^{-120}$; unrounded factors give the exact
$I = 1.098945\times 10^{-120} = 1.0989\times 10^{-120}$ to five
significant figures. (An earlier version displayed
$I = 1.0990\times 10^{-120}$ throughout: the rounded intermediates
above compound to $1.09895$, which rounds up, while the exact value
rounds to $1.0989$ --- last digit corrected at every occurrence,
external review round 73.)
\end{proof}

\subsection{Variational characterisation}\label{sec:variational}

The three regime boundaries of Section~\ref{sec:regimes} are continuous
crossings of $p(d)$ that lie strictly between integers:
\[
d_0^\star \approx 6.2569,\qquad
d_1^\star \approx 19.7308,\qquad
d_2^\star \approx 217.6267.
\]
Each boundary therefore presents an integer pair
$(\lfloor d_i^\star\rfloor,\;\lceil d_i^\star\rceil) =
(6,7),\,(19,20),\,(217,218)$, and the bilinear invariant form of
Theorem~\ref{thm:invariant} takes one of eight possible numerical values
depending on which integer in each pair is substituted for
$\Omega_{d_i}$. No rounding convention selects a canonical integer
uniformly across the three boundaries. The cascade invariant is
defined instead by taking the supremum of the bilinear form over all
eight labelings.

\begin{theorem}[Variational form]\label{thm:variational}
\[
I \;=\; \max_{(d_0,d_1,d_2)}\;
\left(\frac{\Omega_5}{\Omega_{d_0}}\right)^{\!2}
\cdot\;\Omega_{d_1}\cdot\Omega_{d_2},
\qquad
d_0\in\{6,7\},\;
d_1\in\{19,20\},\;
d_2\in\{217,218\}.
\]
The supremum is attained uniquely at $(d_0,d_1,d_2) = (7,19,217)$,
giving $I = 9\,\Omega_{19}\,\Omega_{217}/\pi^2
= 1.0989\times 10^{-120}$.
\end{theorem}

\begin{proof}
The bilinear form is monotone in each integer choice: it is decreasing
in $\Omega_{d_0}$ (which sits in the denominator of the frame correction,
squared) and increasing in $\Omega_{d_1}$ and $\Omega_{d_2}$ (both are
numerator factors). The supremum is therefore attained at the integer
labels that simultaneously minimise $\Omega_{d_0}$ and maximise
$\Omega_{d_1}$ and $\Omega_{d_2}$:
\[
\min(\Omega_6,\Omega_7) = \Omega_7,\quad
\max(\Omega_{19},\Omega_{20}) = \Omega_{19},\quad
\max(\Omega_{217},\Omega_{218}) = \Omega_{217}.
\]
These three extremal choices give $(d_0,d_1,d_2) = (7,19,217)$. Full
enumeration of all eight labelings is given in
\texttt{tools/verifiers/verify\_continuous\_boundary.py}; the argmax is unique
and its value is $1.0989\times 10^{-120}$.
\end{proof}

\begin{remark}[The variational form replaces integer labelling conventions]
\label{rem:variational}
The variational definition eliminates the need to defend a specific
rounding rule. Alternatives such as uniform-floor $(6,19,217)$,
uniform-ceiling $(7,20,218)$, nearest-integer $(6,20,218)$, and the
continuous-crossing value all give smaller invariants: uniform floor
differs by $3.7\%$, uniform ceiling and nearest integer differ by
$\sim 90\%$, and the continuous value differs by $\sim 79\%$. The
supremum $(7,19,217)$ is the only labelling that reproduces the
observation: Planck 2018 gives $\rho_\Lambda/M_{\mathrm{Pl,red}}^4 =
(7.150\pm 0.13)\times 10^{-121}$ (Part~I), which pulled back to
invariant units through Part~I's closure
$\rho_\Lambda/M^4 = (2/\pi)\,e^{0.02108}\,I$ corresponds to
$I_{\mathrm{obs}} = (\pi/2)\,e^{-0.02108}\cdot(7.150\pm 0.13)\times
10^{-121} = (1.10\pm 0.02)\times 10^{-120}$, matched by the supremum
to $0.1\%$ --- well inside the $1.9\%$ Planck $1\sigma$ uncertainty
(Part~I). (An earlier version of this
sentence labelled $(1.10\pm 0.02)\times 10^{-120}$ directly as ``the
observed $\rho_\Lambda/M_{\mathrm{Pl,red}}^4$''; that figure is the
observation in invariant units --- unit label corrected, external
review round 69.) The $0.1\%$ match is therefore the statement that
the supremum of the cascade invariant over boundary labelings
coincides with the observed vacuum energy density expressed in the
invariant's own units, not that a particular
rounding rule happens to hit the right value.
\end{remark}

\begin{remark}[Uniqueness of max over min]
The variational form selects the supremum rather than the infimum.
The infimum $(6,20,218)$ gives $1.02\times 10^{-121}$, an order of
magnitude ($\approx 10.8\times$) below observation. (An earlier
version of this sentence said ``two orders of magnitude''; the ratio
$1.10\times 10^{-120}/1.02\times 10^{-121}\approx 10.8$ is one order
--- corrected, external review round 62.) Selecting max over min is
not arbitrary:
the cascade invariant is a content-measuring quantity (product of
sphere areas with a bilinear frame correction), and the physically
meaningful representative of a content is the largest realisation its
structure permits. A principled derivation of max from the cascade's
own axioms---connecting the supremum to a distinguished quantity such
as an entropy, a boundary area, or a characteristic of the observer's
layer---remains open. (Three exact reformulations narrow the question:
over the eight labelings the supremum is equivalently (i) the
labelling minimising the de~Sitter horizon entropy
$S = 24\pi^2 M^4_{\mathrm{Pl,red}}/\rho_\Lambda$, which is strictly
decreasing in the invariant; (ii) the \emph{odd} member of each
straddling pair---the Euler-null sphere, $\chi(S^d) = 0$, the same
Euler characteristic whose value $\chi(S^{2n}) = 2$ weights the
potential-shift family below (Section~\ref{sec:inter-layer-coupling}
onward); indeed all four distinguished
dimensions $\{5, 7, 19, 217\}$ are odd; and (iii) arithmetically: the
supremum's twist points $s = d+1 \in \{8, 20, 218\}$ are even---the
Euler-rational points $\zeta(s) \in \mathbb{Q}\cdot\pi^s$, whose
functional-equation mirrors $\zeta(1-s)$ are nonzero rationals---while
the infimum's twist points $\{7, 21, 219\}$ mirror exactly onto the
trivial zeros $\zeta(-6) = \zeta(-20) = \zeta(-218) = 0$: the supremum
is the unique labelling avoiding the trivial-zero mirror set. A
principled derivation of any one of these equivalents would settle
the question; none is yet derived.)
\end{remark}

% ============================================================
\section{The Hierarchy}\label{sec:the-hierarchy}
% ============================================================

\begin{theorem}\label{thm:hierarchy}
$\log_{10}(1/I)\approx
\pi e^{2\sqrt{\pi}}(2\sqrt{\pi}-1)/\ln 10\approx 120$.
\end{theorem}

\begin{proof}
Stirling at $d_2 = 2\pi e^{2\sqrt{\pi}}$:
$\ln\Omega_{d_2}\approx -\pi e^{2\sqrt{\pi}}(2\sqrt{\pi}-1)
\approx -277$ nats $\approx -120.3$ decades. The scale correction
$\ln(9/\pi^2) = -0.09$ contributes $0.04$ decades, giving
$\log_{10}(1/I) \approx 120.0$.
\end{proof}

\begin{remark}
$\sqrt{\pi}\approx 1.77$ (orthogonal compression);\;
$e^{2\sqrt{\pi}}\approx 34.6$ (exponential amplification);\;
$2\pi e^{2\sqrt{\pi}}\approx 218$ (the cascade dimension);\;
$(2\sqrt{\pi}-1)\approx 2.54$ (per-dimension suppression).\;
Product: $\approx 277$ nats $\approx 120$ decades.
\end{remark}

% ============================================================
\section{Robustness}\label{sec:robustness}
% ============================================================

The cascade invariant $I = 1.0989\times 10^{-120}$ is a pure number from
the Gamma function. It was not constructed to match any physical quantity.
The identification of $I$ with the dimensionless cosmological constant,
and the observer-frame conversion that makes the numerical comparison
well-defined, are the job of Part~I; this paper records the pure number.
The table below compares alternative cascade invariants to the order-0
target $\log_{10}|I|\approx -120$ set by the four-distinguished-dimension
tower, showing that the full bilinear invariant is the closest among
candidates built from the same primitive set.

\begin{center}
\begin{tabular}{llcl}
\toprule
Invariant & Dimensions & $\log_{10} I$ & Deviation \\
\midrule
Full ($\Omega_5^2\,\Omega_{19}\,\Omega_{217}/\Omega_7^2$)
& 5, 7, 19, 217 & $-119.96$ & $\sim 0.1\%$ \\
Partial ($\Omega_{19}\,\Omega_{217}$)
& 19, 217 & $-119.92$ & $\sim 10\%$ \\
Volume ratio ($V_7/V_5\times\Omega_{19}\,\Omega_{217}$)
& 5, 7, 19, 217 & $-119.97$ & $\sim 1.7\%$ \\
$\{0, c_2\}$ & 7, 217 & $-118.12$ & Factor 68 \\
$\{c_1, 2c_1\}$ & 19, 61 & $-17.3$ & $10^{3}$ OOM \\
$\{2c_1, c_2\}$ & 61, 217 & $-136.6$ & 17 OOM low \\
$\{c_1, c_2\}$, quotient & 19, 217 & $+119.6$ & Wrong sign \\
\bottomrule
\end{tabular}
\end{center}

The full invariant dominates the partial invariant. The volume-ratio
variant confirms that sphere areas, not volumes, are primary
(Corollary~\ref{cor:primary}): using derived quantities instead of
primary quantities costs an order of magnitude in structural proximity
to the target $\log_{10}|I|\approx -120$. No alternative primitive set
comes within an order of magnitude. The ``Deviation'' column measures
the log-space distance from the bilinear target set by the
variational characterisation (Theorem~\ref{thm:variational}), not a
physical observation; numerical comparison with the observed vacuum
energy density is developed in Part~I.

% ============================================================
\section{The Invariant Hierarchy}\label{sec:the-invariant-hierarchy}
% ============================================================

\begin{definition}[Invariant order]\label{def:order}
The \emph{order} of a cascade invariant is the level of geometric
structure it depends on.

\emph{Order~0 (global):} Invariants built from the four distinguished
dimensions alone. The cascade invariant $I$ is the unique order-0
invariant. It depends only on the coarsest features of the Gamma
function: the locations and sphere areas of its equilibria and thresholds.

\emph{Order~1 (structural):} Invariants depending on the layer-by-layer
structure between distinguished dimensions. Examples: the Bott periodicity
classification ($d\bmod 8$), the layer count $d_2 - d_1 = 198$, the
boundary dominance ratio $\Omega_{d-1}/V_d = d$, individual recurrence
ratios $\Omega_{d+1}/\Omega_d$.

\emph{Order~2 (propagation):} Invariants depending on the accumulated
cascade geometry across multiple layers. Examples: the cascade potential
$\Phi(d) = \sum p(k)$, the lapse function $N(d) = \sqrt{\pi}\cdot R(d)$,
fermion mass ratios from multi-step descent.
\end{definition}

\begin{remark}[Order and precision]
The order-0 invariant $I$ is built from the four distinguished dimensions
alone and is immune to the inter-layer coupling that generates the
subleading corrections appearing in the order-1 and order-2 predictions
of the companion papers. Whether any given order is quantitatively close
to a physical observable is a question for the companion papers, where
the physical identification hypothesis is introduced; Part~0 records only
the structural distinction that the order-0 invariant does not accumulate
the descent-dependent corrections that affect the higher orders. The
inter-layer coupling itself is computed in the next section, completing
the structural account of how an observable's traversal of multiple
cascade layers departs from the independent-step approximation.
\end{remark}

% ============================================================
\section{Inter-Layer Coupling}\label{sec:inter-layer-coupling}
% ============================================================

The slicing recurrence treats each compactification step as independent: the propagator through $n$ steps is the product of $n$ individual lapses. This is exact for volumes (the telescoping identity $V_D/V_4 = \prod N(d)$ holds by the Gamma function). It is not exact for quantities that depend on the overlap structure of consecutive layer integrands. The correction is determined by the eigenvalue deficit of the Gram correlation matrix of the integrands $f_d(x) = (1 - x^2)^{d/2}$---a Beta function identity with no free parameters.

The Gram correlation has an exact closed form in cascade-native primitives (Theorem~\ref{thm:gram-R}): $C^2_{d,d+1} = R(2d+2)^2 / [R(2d+1)\,R(2d+3)]$, equivalently $\log C^2_{d,d+1} = -\tfrac{1}{2}\Delta^2 \log\alpha\big|_{2d+2}$ (Corollary~\ref{cor:gram-laplacian}). The Gram correction is therefore the discrete log-Laplacian of the cascade compliance function $\alpha(d) = R(d)^2/4$ (Definition~\ref{def:compliance-action}), evaluated at the doubled cascade argument. The same compliance function admits a second cascade-internal derivative --- the marginal Green's function $G(d^*) - G(d^*+1) = \alpha(d^*)$ (Proposition~\ref{prop:marginal-greens}) --- which is the linear source response used downstream by Part~IVb's $\alpha(d^*)/\chi^k$ precision-closure family. The path-distributed correction and that single-source family are therefore complementary derivatives of one cascade-internal object, both established within Part~0.

Throughout this section: ``second order'' refers exclusively to the perturbation-theory order in the eigenvalue expansion (Theorem~\ref{thm:pert}).

\subsection{The Gram matrix of cascade layer integrands}

\begin{definition}[Gram matrix]
For a set of cascade layers $\{d_1, d_2, \ldots, d_n\}$, the Gram matrix is
\[
G_{ij} = \langle f_{d_i}, f_{d_j} \rangle = \int_{-1}^{1} (1 - x^2)^{(d_i + d_j)/2}\, dx.
\]
\end{definition}

\begin{theorem}[Gram matrix as Beta function]\label{thm:gram}
\[
G_{ij} = B\!\left(\tfrac{1}{2},\, \tfrac{d_i + d_j}{2} + 1\right) = \sqrt{\pi}\, \frac{\Gamma\!\left(\frac{d_i + d_j}{2} + 1\right)}{\Gamma\!\left(\frac{d_i + d_j}{2} + \frac{3}{2}\right)}.
\]
\end{theorem}

\begin{proof}
Substituting $t = x^2$ in the integral: $\int_{-1}^{1}(1 - x^2)^\alpha\, dx = B(1/2, \alpha + 1)$ with $\alpha = (d_i + d_j)/2$. The Beta function identity $B(a,b) = \Gamma(a)\Gamma(b)/\Gamma(a+b)$ with $a = 1/2$ gives the stated form.
\end{proof}

\begin{definition}[Correlation matrix]
The normalised Gram matrix (correlation matrix) is
\[
C_{ij} = \frac{G_{ij}}{\sqrt{G_{ii}\, G_{jj}}} = \frac{B\!\left(\frac{1}{2},\, \frac{d_i + d_j}{2} + 1\right)}{\sqrt{B\!\left(\frac{1}{2},\, d_i + 1\right)\, B\!\left(\frac{1}{2},\, d_j + 1\right)}}.
\]
$C_{ii} = 1$ for all $i$, $\mathrm{tr}\, C = n$, and $0 < C_{ij} < 1$ for $i \neq j$.
\end{definition}

\subsection{The adjacent-layer overlap deficit}

\begin{theorem}[Overlap deficit]\label{thm:overlap}
For adjacent layers $d$ and $d+1$:
\[
1 - C_{d,d+1}^2 = 1 - \frac{B\!\left(\frac{1}{2},\, d + \frac{3}{2}\right)^2}{B\!\left(\frac{1}{2},\, d+1\right)\, B\!\left(\frac{1}{2},\, d+2\right)}.
\]
This quantity is strictly positive for all $d \geq 1$ and satisfies:
\[
1 - C_{d,d+1}^2 \sim \frac{1}{8\, d^2} \quad (d \to \infty).
\]
Numerically: $1 - C_{5,6}^2 = 0.00319$, $\; 1 - C_{11,12}^2 = 0.00083$, $\; 1 - C_{19,20}^2 = 0.00030$.
\end{theorem}

\begin{proof}
\emph{Positivity.} $C_{d,d+1} < 1$ by Cauchy--Schwarz: the integrands $f_d$ and $f_{d+1}$ are not proportional in $L^2[-1,1]$ for $d \neq d+1$, so the normalised inner product is strictly less than~$1$.

\emph{Asymptotic $1 - C_{d,d+1}^2 \sim 1/(8d^2)$.} Set $h(\alpha) = B(1/2,\alpha+1) = \sqrt{\pi}\,\Gamma(\alpha+1)/\Gamma(\alpha+3/2)$ and $L(\alpha) = \log h(\alpha)$. Then
\[
\log C_{d,d+1}^2 \;=\; 2L(d+\tfrac{1}{2}) - L(d) - L(d+1),
\]
the centered second difference of $L$ at $d+1/2$ with step $1/2$. Taylor expanding,
\[
L(d) + L(d+1) \;=\; 2L(d+\tfrac{1}{2}) + \tfrac{1}{4}L''(d+\tfrac{1}{2}) + O(L^{(4)}/192),
\]
so $\log C_{d,d+1}^2 = -\tfrac{1}{4}L''(d+\tfrac{1}{2}) + O(L^{(4)})$.

Now $L'(\alpha) = \psi(\alpha+1) - \psi(\alpha+3/2)$, where $\psi$ is the digamma function. The asymptotic expansion $\psi(x) = \log x - 1/(2x) - 1/(12x^2) + O(1/x^4)$ gives, after expanding $\psi(\alpha+1)$ and $\psi(\alpha+3/2)$ around $\alpha$ and combining,
\[
L'(\alpha) \;=\; -\frac{1}{2\alpha} + \frac{3}{8\alpha^2} + O(1/\alpha^3),
\qquad
L''(\alpha) \;=\; \frac{1}{2\alpha^2} + O(1/\alpha^3).
\]
Therefore
\[
\log C_{d,d+1}^2 \;=\; -\frac{1}{4}\cdot\frac{1}{2d^2} + O(1/d^3) \;=\; -\frac{1}{8d^2} + O(1/d^3),
\]
and $1 - C_{d,d+1}^2 = -\log C_{d,d+1}^2 + O\bigl((\log C^2)^2\bigr) = 1/(8d^2) + O(1/d^3)$.

The numerical values quoted in the theorem statement ($1 - C_{5,6}^2 = 0.00319$, $1 - C_{11,12}^2 = 0.00083$, $1 - C_{19,20}^2 = 0.00030$) are evaluated directly from the Beta function ratio; the asymptotic $1/(8d^2)$ supplies the leading $1/d^2$ scaling and the sign, but subleading $1/d^3$ terms matter at small $d$ (the asymptotic over-estimates by $\sim 50\%$ at $d=5$, by $\sim 25\%$ at $d=11$, and by $\sim 15\%$ at $d=19$, converging to the exact value as $d\to\infty$).
\end{proof}

\begin{remark}
The overlap deficit $1 - C_{d,d+1}^2$ is the geometric content at step $d$ that is \emph{not} shared with step $d+1$. It measures the imperfect folding: the fraction by which consecutive layer integrands fail to be collinear in $L^2$. Its positivity is a consequence of $R_{\mathrm{eff}}(d) = 1/\sqrt{d+3} > 0$ (Theorem~\ref{thm:reff}): the compactification at each step is asymptotic, never complete.
\end{remark}

\subsection{Closed form via cascade slicing ratios}\label{sec:gram-closed-form}

The asymptotic of Theorem~\ref{thm:overlap} extends to an \emph{exact} closed form expressing $C^2_{d,d+1}$ in cascade-native primitives. The exact form reveals that the Gram correlation is the discrete log-Laplacian of the cascade scalar action's compliance function $\alpha(d) = R(d)^2/4$, evaluated at the doubled cascade argument $2d+2$.

\begin{theorem}[Gram correlation in cascade slicing ratios]\label{thm:gram-R}
For all $d \geq 1$, the Gram correlation between adjacent cascade layers is
\begin{equation}\label{eq:gram-R-form}
C^2_{d,d+1} \;=\; \frac{R(2d+2)^2}{R(2d+1)\,R(2d+3)},
\end{equation}
where $R(d) = \Gamma((d+1)/2)/\Gamma((d+2)/2)$ is the cascade slicing ratio of \S\ref{sec:recurrence}.
\end{theorem}

\begin{proof}
The Gram entries (Theorem~\ref{thm:gram}) are Beta functions $G_{ij} = B(1/2, (d_i+d_j)/2 + 1)$. Using the Beta--Gamma reduction $B(1/2, n+1) = \sqrt{\pi}\,\Gamma(n+1)/\Gamma(n+3/2)$ together with the cascade slicing ratio's defining identity $R(d) = \Gamma((d+1)/2)/\Gamma((d+2)/2)$, one obtains:
\[
G_{d,d}     = B(\tfrac{1}{2}, d+1)   = \sqrt{\pi}\,R(2d+1), \qquad
G_{d,d+1}   = B(\tfrac{1}{2}, d+\tfrac{3}{2}) = \sqrt{\pi}\,R(2d+2),
\]
\[
G_{d+1,d+1} = B(\tfrac{1}{2}, d+2)   = \sqrt{\pi}\,R(2d+3).
\]
Substituting into the definition $C^2_{d,d+1} = G_{d,d+1}^2 / (G_{d,d}\,G_{d+1,d+1})$ gives
$C^2_{d,d+1} = R(2d+2)^2 / [R(2d+1)\,R(2d+3)]$.
\end{proof}

\begin{definition}[Cascade compliance and scalar action]\label{def:compliance-action}
The \emph{cascade compliance} at layer $d$ is
\[
\alpha(d) \;:=\; \frac{R(d)^2}{4},
\]
the square of the cascade slicing ratio of \S\ref{sec:recurrence}, normalised by 4. As a $\Gamma$-function ratio, $\alpha$ is a cascade-internal quantity with no free parameter.

The \emph{cascade scalar action} on the layer index is the discrete elastic functional
\[
S[\varphi] \;:=\; \sum_d \frac{1}{2\alpha(d)}\bigl(\Delta\varphi(d)\bigr)^2,
\qquad \Delta\varphi(d) := \varphi(d+1) - \varphi(d).
\]
By construction, $S$ is the unique quadratic nearest-neighbour functional whose Euler--Lagrange equation has the cascade slicing recurrence as equilibrium; the compliance $\alpha(d) = R(d)^2/4$ is the cascade-natural choice making the action's marginal Green's function equal $\alpha(d^*)$ at every distinguished source layer $d^*$ (Proposition~\ref{prop:marginal-greens} below).
\end{definition}

\begin{proposition}[Marginal Green's function of the cascade scalar action]\label{prop:marginal-greens}
The cumulative compliance from a cascade source layer $d^*$ to the cascade terminus $d_2 = 217$ is
\[
G(d^*) \;:=\; \sum_{k = d^*}^{d_2 - 1} \alpha(k),
\]
and the per-layer marginal contribution at source $d^*$ is the local compliance value:
\[
G(d^*) - G(d^* + 1) \;=\; \alpha(d^*).
\]
This is the source-response identity used by Part~IVb's $\alpha(d^*)/\chi^k$ precision-closure family.
\end{proposition}

\begin{proof}
Direct: $G(d^*)$ and $G(d^* + 1)$ differ by exactly the term at $k = d^*$, which is $\alpha(d^*)$.
\end{proof}

\begin{corollary}[Gram-Laplacian identity]\label{cor:gram-laplacian}
The log-correlation between adjacent cascade layers is the discrete log-Laplacian of the cascade compliance $\alpha$ (Definition~\ref{def:compliance-action}), evaluated at the doubled argument $2d+2$ with unit step:
\begin{equation}\label{eq:gram-laplacian}
\log C^2_{d,d+1} \;=\; -\frac{1}{2}\,\Delta^2 \log\alpha\,\Big|_{2d+2},
\end{equation}
where $\Delta^2 f|_n := f(n-1) + f(n+1) - 2f(n)$ is the centered discrete Laplacian.
\end{corollary}

\begin{proof}
By Theorem~\ref{thm:gram-R} and properties of the logarithm,
\[
\log C^2_{d,d+1} = 2\log R(2d+2) - \log R(2d+1) - \log R(2d+3) = -\Delta^2 \log R\,\big|_{2d+2}.
\]
Since $\alpha(d) = R(d)^2/4$, we have $\log\alpha = 2\log R - 2\log 2$, so $\Delta^2 \log\alpha = 2\,\Delta^2 \log R$ (constants drop in second differences). Substituting yields equation~\eqref{eq:gram-laplacian}.
\end{proof}

\begin{remark}[Numerical verification]\label{rem:gram-verifier}
Both Theorem~\ref{thm:gram-R} (the closed form) and Corollary~\ref{cor:gram-laplacian} (the Laplacian identity) are verified at all tested $d$ values in
\begin{center}
\texttt{tools/verifiers/gram\_compliance\_laplacian.py}
\end{center}
to machine precision on the closed form and to $\sim 10^{-7}$ relative on the Laplacian identity (the floating-point cancellation noise floor at large $d$, where $\log C^2_{d,d+1} \sim 1/(8d^2)$ is computed as a difference of $O(1)$ log-Gamma terms).
\end{remark}

\begin{remark}[Structural unification with the cascade scalar action]\label{rem:gram-action-unification}
Corollary~\ref{cor:gram-laplacian} establishes a direct structural connection between the Gram framework above and the cascade scalar action $S[\varphi]$ of Definition~\ref{def:compliance-action}. Both corrections derivable from $S[\varphi]$ and its compliance $\alpha(d)$ have a cascade-internal expression:
\begin{itemize}
\item \emph{Linear source response} (single-source, used downstream by Part~IVb): the action's marginal Green's function at source layer $d^*$ equals $\alpha(d^*)$ exactly (Proposition~\ref{prop:marginal-greens}). Chirality factorisation through $k$ independent channels gives the $1/\chi^k$ attenuation that produces the cascade's $\alpha(d^*)/\chi^k$ precision-closure family.
\item \emph{Discrete log-Laplacian of the compliance} (path-distributed, this section): the per-step Gram correction is $\frac{1}{2}\Delta^2(-\log\alpha)|_{2d+2}$ (Corollary~\ref{cor:gram-laplacian}). It is positive because $-\log\alpha$ is convex on the cascade tower, i.e.\ $\alpha$ is log-concave.
\end{itemize}
The two correction families are complementary in the action's dual variables: zeroth-order source response (the marginal $\alpha(d^*)$ of a single layer) versus second-order log-compliance curvature (the Laplacian of $\log\alpha$ summed across a path). Both are derived within Part~0 from the same compliance function. The doubling $d\mapsto 2d+2$ is a Beta--Gamma reduction artefact: cascade integrand inner products $\langle f_d, f_{d'}\rangle$ produce $R$-values at doubled arguments via $B(1/2, n+1) = \sqrt{\pi}\,R(2n+1)$ and $B(1/2, n+3/2) = \sqrt{\pi}\,R(2n+2)$, so the Gram correlation between layers $d$ and $d+1$ depends on cascade $R$-values at doubled indices $\{2d+1, 2d+2, 2d+3\}$.
\end{remark}

\subsection{The eigenvalue structure of the correlation matrix}

\begin{theorem}[Near-collinearity]\label{thm:eigenvalue}
For a path of $n$ consecutive layers $\{d_0, d_0+1, \ldots, d_0+n-1\}$, the correlation matrix $C$ has one dominant eigenvalue $\lambda_1$ satisfying
\[
\frac{\lambda_1}{n} = 1 - \epsilon(n, d_0), \qquad \epsilon > 0,
\]
and $(n - 1)$ subdominant eigenvalues summing to $n\,\epsilon$. The deficit $\epsilon$ depends on both $n$ (the number of layers) and $d_0$ (where in the cascade the path begins).
\end{theorem}

\begin{proof}
The entries $C_{ij}$ are smooth functions of $|d_i - d_j|$ that decrease from $C_{ii} = 1$ towards a minimum at $|d_i - d_j| = n - 1$. For the path $d = 5, \ldots, 12$ ($n = 8$): the minimum entry is $C_{5,12} = 0.9622$, far from zero. The matrix is therefore close to the rank-1 matrix $\mathbf{1}\mathbf{1}^T/n$, whose sole nonzero eigenvalue is $n$. The deficit $\epsilon$ is determined by the off-diagonal decay rate of $C_{ij}$, which is controlled by the Beta function asymptotics of Theorem~\ref{thm:overlap}.
\end{proof}

\noindent Numerical eigenvalue structure for $d = 5, \ldots, 12$ ($n = 8$):

\begin{center}
\renewcommand{\arraystretch}{1.2}
\begin{tabular}{ccc}
\toprule
Eigenvalue & Value & Fraction of trace \\
\midrule
$\lambda_1$ & 7.9347263 & 99.184\% \\
$\lambda_2$ & 0.0646528 & 0.808\% \\
$\lambda_3$ & 0.0006152 & 0.008\% \\
$\lambda_4$--$\lambda_8$ & $< 10^{-4}$ & $< 0.001\%$ \\
\midrule
$\epsilon = 1 - \lambda_1/n$ & \multicolumn{2}{c}{$0.008159 = 0.816\%$} \\
\bottomrule
\end{tabular}
\end{center}

\noindent The cascade layers are 99.2\% collinear. The remaining 0.8\% is the inter-layer coupling.

\subsection{Path dependence of the eigenvalue deficit}

The deficit $\epsilon(n, d_0)$ depends on where in the cascade the path lies. Lower $d$ (closer to the observer, larger $R_{\mathrm{eff}}$, less complete folding) gives larger deficits.

\begin{center}
\renewcommand{\arraystretch}{1.2}
\begin{tabular}{llcc}
\toprule
Path & $n$ & $\epsilon = 1 - \lambda_1/n$ & Needed correction \\
\midrule
$d = 5, \ldots, 12$ & 8 & 0.00816 & $1.73\%$ ($\alpha_s$) \\
$d = 6, \ldots, 13$ & 8 & 0.00657 & $1.75\%$ ($m_\tau/m_\mu$) \\
$d = 10, \ldots, 17$ & 8 & 0.00332 & --- \\
$d = 14, \ldots, 21$ & 8 & 0.00200 & $0.13\%$ ($m_\mu/m_e$) \\
\bottomrule
\end{tabular}
\end{center}

\noindent At fixed $n$, $\epsilon$ decreases with $d_0$ because the overlap deficit per step (Theorem~\ref{thm:overlap}) decreases as $1/d^2$.

\subsection{The perturbation theory for \texorpdfstring{$\epsilon$}{epsilon}}\label{sec:perturbation}

\begin{theorem}[First-order eigenvalue deficit]\label{thm:pert}
To first order in the deficit matrix $D = J - C$,
\[
\epsilon \;=\; 1 - \frac{\lambda_1}{n} \;=\; \frac{2}{n^2} \sum_{i < j} \left(1 - C_{ij}\right) \;+\; O\!\bigl(\|D\|^2\bigr),
\]
where $C_{ij}$ are the entries of the correlation matrix. The second-order correction is \emph{positive} (a sum of squared matrix elements of $D$ in the zero-eigenspace of $J$; see proof), so the right-hand side is an \emph{upper bound} for the eigenvalue deficit $\epsilon$.
\end{theorem}

\begin{proof}
Write $C = J - D$ where $J_{ij} = 1$ for all $i,j$ and $D_{ij} = 1 - C_{ij} \geq 0$, with $D_{ii} = 0$. The matrix $J$ has the non-degenerate eigenvalue $n$ with eigenvector $u = (1, \ldots, 1)/\sqrt{n}$, and eigenvalue $0$ with multiplicity $n - 1$ (eigenspace $u^\perp$). By standard Rayleigh--Schr\"odinger perturbation theory for the non-degenerate eigenvalue $n$:
\[
\lambda_1(C) \;=\; n - \langle u | D | u \rangle \;+\; \frac{1}{n}\sum_{v \in u^\perp}|\langle v | D | u \rangle|^2 \;+\; O\!\bigl(\|D\|^3\bigr).
\]
The first-order term evaluates to
\[
\langle u | D | u \rangle \;=\; \frac{1}{n}\sum_{i,j}D_{ij} \;=\; \frac{2}{n}\sum_{i<j}(1 - C_{ij}),
\]
giving the stated first-order formula after dividing by $n$. The second-order term is a sum of squared matrix elements, hence non-negative; it vanishes only when $D$ has no component in the $u^\perp$ sector (the rank-1 case, when $C$ is a constant-off-diagonal matrix). In all non-trivial cases it is strictly positive, so $\lambda_1(C) > n - \langle u|D|u\rangle$ at second order, equivalently $\epsilon < (2/n^2)\sum_{i<j}(1 - C_{ij})$: the first-order formula over-estimates the deficit, and the right-hand side is an upper bound.
\end{proof}

\noindent Verification: for $d = 5, \ldots, 12$ ($n = 8$), the first-order formula gives $\epsilon^{(1)} = 0.008173$; the exact eigenvalue gives $\epsilon = 0.008159$. Agreement: $0.17\%$; the exact value is below the first-order formula as predicted by the upper-bound direction of the theorem. The formula is verified to comparable precision for all paths tested.

\begin{corollary}[Closed form for $\lambda_1$ at $n = 2$]\label{cor:lambda1-n2}
For a two-layer path $\{d_0, d_0 + 1\}$, the dominant eigenvalue admits the exact closed form
\[
\lambda_1(2, d_0) \;=\; 1 + \frac{R(2 d_0 + 2)}{\sqrt{R(2 d_0 + 1)\,R(2 d_0 + 3)}}
\;=\; 1 + \exp\!\Bigl(-\tfrac{1}{4}\,\Delta^2 \log\alpha\big|_{2d_0+2}\Bigr),
\]
where the second equality follows from Corollary~\ref{cor:gram-laplacian}. Equivalently, $\epsilon(2, d_0) = 1 - \lambda_1/2 = (1 - C_{d_0, d_0+1})/2$, and the perturbation-theory upper bound of Theorem~\ref{thm:pert} is saturated at $n = 2$.
\end{corollary}

\begin{proof}
At $n = 2$ the matrix $C = \begin{pmatrix}1 & c \\ c & 1\end{pmatrix}$ with $c = C_{d_0, d_0 + 1}$ has eigenvectors $u_\pm = (1, \pm 1)/\sqrt{2}$ and eigenvalues $1 \pm c$, so $\lambda_1 = 1 + c$. The uniform vector $u = u_+$ is exactly the dominant eigenvector; the second-order correction in the proof of Theorem~\ref{thm:pert} vanishes because $u^\perp$ is one-dimensional and orthogonal to $u$ by the symmetry of $C$. Substituting Theorem~\ref{thm:gram-R} (and its Laplacian form, Corollary~\ref{cor:gram-laplacian}) for $c$ gives the stated closed form.
\end{proof}

\begin{remark}[Status of the $\lambda_1$ closed form for $n \geq 3$]\label{rem:lambda1-no-closed-form}
For $n \geq 3$ no single-equation radical form for $\lambda_1$ in cascade slicing ratios has been identified. What is available is the power-trace bound family $U_k = (\mathrm{tr}\,C^k)^{1/k}$ (Corollary~\ref{cor:lambda1-power-trace}), which provides closed-form approximations to $\lambda_1$ converging at geometric rate; combined with the $n = 2$ exact identity (Corollary~\ref{cor:lambda1-n2}) and the perturbation-theory bound of Theorem~\ref{thm:pert}, these cover all cascade-physics paths to better than machine precision.
\end{remark}

\begin{remark}[OP factorisation of the cascade correlation matrix]\label{rem:lambda1-OP-factorisation}
$C$ factorises as $C = A A^{\top}$ with $A$ lower triangular and $A_{ik} = \langle p_k, x^i\rangle_{\mu_m} / \|x^i\|_{\mu_m}$, where $\{p_k\}$ are the orthonormal polynomials of the measure $d\mu_m(x) = x^m/\sqrt{1-x^2}\,dx$ on $[0,1]$ ($m = 2 d_0 + 1$). Hence $\lambda_k(C) = \sigma_k(A)^2$. Under $u = x^2$, $\mu_m$ is a shifted Jacobi measure, identifying the closed-form question for $\lambda_1$ as a question in classical orthogonal-polynomial theory (\texttt{tools/research/lambda1\_jacobi\_route.py}).
\end{remark}

\begin{corollary}[Power-trace bound family closing $\lambda_1$ to arbitrary precision]\label{cor:lambda1-power-trace}
For each $k \geq 1$,
\[
U_k \;=\; \bigl(\mathrm{tr}\,C^k\bigr)^{1/k} \;\geq\; \lambda_1,
\]
with $U_1 \geq U_2 \geq \cdots \geq \lambda_1$ monotone, and $U_k \to \lambda_1$ at geometric rate $(\lambda_2/\lambda_1)^k$. Each $U_k$ is closed-form in cascade slicing ratios:
\[
\mathrm{tr}\,C^k \;=\; \sum_{i_1, \ldots, i_k} C_{i_1 i_2}\,C_{i_2 i_3} \cdots C_{i_k i_1},
\]
a polynomial sum of $R$-ratio products. Numerically, for cascade-physics paths ($d_0 \in \{5, \ldots, 19\}$, $n \leq 10$): $U_2$ is tight to $\sim 10^{-5}$, $U_3$ to $\sim 10^{-9}$, $U_4$ to $\sim 10^{-12}$, $U_6$ to machine precision; for the long $\rho_\Lambda$ descent ($n = 26$), $U_4$ is tight to $\sim 10^{-7}$ (\texttt{tools/research/lambda1\_power\_trace\_bounds.py}). The bound family thus determines $\lambda_1$ in closed form to any desired precision.
\end{corollary}

\begin{proof}
For positive eigenvalues $\lambda_1 \geq \lambda_2 \geq \cdots \geq 0$, $\mathrm{tr}\,C^k = \sum_j \lambda_j^k$, so $\lambda_1^k \leq \mathrm{tr}\,C^k$ and $\lambda_1 \leq U_k$. Monotonicity $U_k \geq U_{k+1}$ is the standard fact that $\ell^k$ norms of a positive vector decrease in $k$. The convergence $U_k \to \lambda_1$ at rate $(\lambda_2/\lambda_1)^k$ follows from $\mathrm{tr}\,C^k = \lambda_1^k(1 + (\lambda_2/\lambda_1)^k + \cdots)$ via the Bernoulli expansion of $(1 + \epsilon)^{1/k}$. The closed-form expression for $\mathrm{tr}\,C^k$ is the standard polynomial trace expansion; each $C_{ij}$ is a closed-form ratio of slicing ratios (Theorem~\ref{thm:gram-R} and its generalisation to arbitrary index pairs), so the sum is closed-form in cascade primitives.
\end{proof}

\subsection{The independent-step correction}\label{sec:first-order}

The cascade's leading-order predictions treat each compactification step as independent. A quantity that depends on the cascade's geometry over a path of $n$ layers receives a correction from the inter-layer coupling. Two distinct measures of this coupling are derived from the Gram matrix:

\begin{itemize}
\item \emph{Eigenvalue deficit $\epsilon = 1 - \lambda_1/n$} (Theorem~\ref{thm:eigenvalue}): the matrix-level magnitude of inter-layer coupling, characterised by the dominant eigenmode of $C$ and aggregated over all $n^2$ pair correlations with weight $(1-C_{ij})$. Derived perturbatively in Section~\ref{sec:perturbation}.
\item \emph{Geometric correction $\delta\Phi$} (Definition~\ref{def:delta-phi} below): the multiplicative log-decrement a propagator picks up traversing the path, summed over adjacent-layer overlap deficits with weight $(1-C^2_{d,d+1})$. Derived in Section~\ref{sec:first-order}.
\end{itemize}

The two are related but structurally distinct: $\epsilon$ characterises the Gram matrix as a whole; $\delta\Phi$ is the cumulative shift along a specific path. The corrected observable formula uses $\delta\Phi$.

\begin{definition}[Geometric correction along a path]\label{def:delta-phi}
For a cascade observable traversing the path $\{d_0, d_0+1, \ldots, d_0+n-1\}$,
\[
\delta\Phi \;:=\; \sum_{k=0}^{n-2}\bigl(1 - C^2_{d_k, d_{k+1}}\bigr),
\]
the path-cumulative adjacent overlap deficit. By Theorem~\ref{thm:overlap}, every term is strictly positive, so $\delta\Phi > 0$ for every cascade path. Theorem~\ref{thm:first-order} below shows that $\delta\Phi$ equals the first-order multiplicative correction to a cascade-descent observable, and Corollary~\ref{cor:exp-resum} shows that the unique closed-form correction respecting path-splitting composition is $Q = Q_0 \cdot \exp(\delta\Phi)$.
\end{definition}

\begin{theorem}[Adjacent overlap deficit correction]\label{thm:first-order}
For a descent-dependent observable traversing path $\{d_0,\ldots,d_0+n-1\}$:
\begin{equation}\label{eq:first-order}
\frac{\delta Q}{Q_0} = \sum_{k=0}^{n-2}\bigl(1 - C_{d_k,d_{k+1}}^2\bigr) \;=\; \delta\Phi,
\end{equation}
where each $C_{d,d+1}$ is the adjacent-layer correlation (Theorem~\ref{thm:overlap}).
This is a sum of Beta function ratios. No free parameter enters.
\end{theorem}

\begin{proof}
The observable propagates through $n$ layers. The independent-step approximation treats each layer as perfectly collinear with the next ($C_{d,d+1}=1$). The first-order departure from collinearity at each step is the squared overlap deficit $1-C_{d,d+1}^2$, which measures the fraction of the integrand at step $d$ that is orthogonal to step $d+1$ (Theorem~\ref{thm:overlap}). For a multiplicative propagator, the leading correction is the sum of per-step deficits, which is $\delta\Phi$ by Definition~\ref{def:delta-phi}.
\end{proof}

\begin{corollary}[Exponential resummation forced by multiplicative composition]\label{cor:exp-resum}
For a multiplicative propagator $Q_0 = \prod_k (\text{step}_k)$
traversing the path of Theorem~\ref{thm:first-order}, the unique
closed-form leading-order Gram-corrected observable respecting
multiplicative composition under path splitting is
\begin{equation}\label{eq:exp-form}
Q \;=\; Q_0\cdot\exp(\delta\Phi) \;+\; O(\delta^2),
\end{equation}
where $\delta = \max_k (1 - C_{d_k,d_{k+1}}^2)$. The exponential
matches Theorem~\ref{thm:first-order} at first order in $\delta$ per
step.
\end{corollary}

\begin{proof}
\emph{Setup.} Because the cascade recurrence is first-order
multiplicative (Corollary~\ref{cor:primary}), the Gram-corrected step at
layer $d_k$ receives a multiplicative factor
$1 + (1 - C_{d_k,d_{k+1}}^2) + O((1-C^2)^2)$, accumulating as a product:
\[
Q \;=\; Q_0\prod_{k=0}^{n-2}\bigl(1 + \delta_k + O(\delta_k^2)\bigr),
\qquad \delta_k := 1 - C_{d_k,d_{k+1}}^2.
\]
This product is the exact corrected observable (modulo $O(\delta^2)$
per-step residuals from Theorem~\ref{thm:first-order}'s first-order
truncation). Define the path-functional
$\Phi(P) := \sum_{k\in P} \delta_k = \delta\Phi(P)$ and the correction factor
$K(P) := Q/Q_0$ over a path $P$.

\emph{Composition law.} Multiplicative composition under path
splitting is a structural property of the cascade recurrence: for a
path $P = P_1 \cup P_2$ formed by concatenating two consecutive
sub-paths, the corrected observable composes as
$Q(P) = Q_0(P_1)\cdot Q_0(P_2) \cdot K(P_1) \cdot K(P_2) = Q_0(P)\cdot K(P_1)\cdot K(P_2)$.
Therefore $K(P) = K(P_1)\cdot K(P_2)$, and since $\Phi(P) = \Phi(P_1) + \Phi(P_2)$
(the path-functional is additive under concatenation), $K$ as a
function of $\Phi$ must satisfy
\begin{equation}\label{eq:cauchy}
K(\Phi_1 + \Phi_2) \;=\; K(\Phi_1) \cdot K(\Phi_2).
\end{equation}

\emph{Cauchy's exponential functional equation forces the form.}
Equation~\eqref{eq:cauchy} together with continuity (a closed-form
correction is a smooth function of its path-functional argument) and
the boundary condition $K(0) = 1$ (no correction for an empty path)
has, by the standard solution of Cauchy's exponential functional
equation, the unique family of solutions $K(\Phi) = e^{c\Phi}$ for
some constant $c$. Matching the first-order behaviour to
Theorem~\ref{thm:first-order} ($K(\Phi) = 1 + \Phi + O(\Phi^2)$ for
small $\Phi$, hence $K'(0) = 1$) fixes $c = 1$. Therefore
$K(\delta\Phi) = e^{\delta\Phi}$, and the corrected observable is~\eqref{eq:exp-form}.

\emph{Why the linear form fails.} The linear closed-form
$K_{\mathrm{lin}}(\Phi) = 1 + \Phi$ matches at first order but
violates~\eqref{eq:cauchy}: $K_{\mathrm{lin}}(\Phi_1) K_{\mathrm{lin}}(\Phi_2)
= (1 + \Phi_1)(1 + \Phi_2) = 1 + \Phi_1 + \Phi_2 + \Phi_1\Phi_2 \neq
1 + (\Phi_1 + \Phi_2) = K_{\mathrm{lin}}(\Phi_1 + \Phi_2)$. The cross-term
$\Phi_1\Phi_2$ breaks additivity. The exponential is the unique
continuous closed form respecting composition; the linear form is the
unique closed form failing it.

\emph{Residual versus the exact product.} The exact corrected
observable is the product $Q_0\prod(1 + \delta_k + O(\delta_k^2))$, not
the exponential. The exponential and the product agree at first order
in $\delta$ per step but differ at second order: the exact product
expands as $1 + \sum\delta_k + \sum_{j<k}\delta_j\delta_k + O(\delta^3)$,
while the exponential expands as $1 + \sum\delta_k + \tfrac{1}{2}(\sum\delta_k)^2 + O(\delta^3)
= 1 + \sum\delta_k + \tfrac{1}{2}\sum\delta_k^2 + \sum_{j<k}\delta_j\delta_k + O(\delta^3)$.
The exponential overestimates the exact product by
$\tfrac{1}{2}\sum\delta_k^2$ at second order. The linear closed-form
$1 + \sum\delta_k$, by contrast, underestimates the product by
$\sum_{j<k}\delta_j\delta_k$ at second order.

For cascade paths starting at $d_0\geq 5$ the per-step deficits scale
as $\delta_k = 1 - C^2_{d_k,d_{k+1}} \sim 1/(8d^2)$
(Theorem~\ref{thm:overlap}), so
$\sum_{k}\delta_k^2 \lesssim \tfrac{1}{64}\sum_{d\geq 5} d^{-4} \approx
6\times 10^{-5}$, and the exponential overshoots the exact product by
at most $\tfrac{1}{2}(6\times 10^{-5}) \approx 3\times 10^{-5}$ in the
multiplicative correction (i.e.\ $\sim 0.003\%$ of $Q$). The linear
form underestimates the product by approximately
$\tfrac{1}{2}(\sum\delta)^2 \approx 2\times 10^{-4}$ ($\sim 0.022\%$
of $Q$) for the longest path ($d=5$ to $d=217$). The exponential is
therefore not only the unique composition-respecting closed form but
also the closer numerical approximation to the exact product than the
linear form, by an order of magnitude.

Both gaps are below the reported observational precisions of all
currently-tested cascade observables ($\sim 0.1\%$ or coarser). The
exponential is therefore the cascade's natural closed-form correction
--- forced by composition, exact at first order in $\delta$, and within
observational precision of the exact product for all paths considered.
It is the form used downstream (Part~I, $\rho_\Lambda$ closure;
Part~IVb, the $v$ and $m_\tau$ absolute-mass shifts).
\end{proof}

\begin{remark}[What the correction is not]
The correction is not to the Gamma function, which is exact. It is not to the slicing recurrence, which is exact (the telescoping identity $V_D/V_4 = \prod N(d)$ holds by $\Gamma(a+1) = a\Gamma(a)$). It is to the \emph{independent-step approximation}: the assumption that a quantity depending on $n$ cascade layers is the product of $n$ individually exact factors. The layers are individually exact but not independent---their integrands overlap in $L^2$, and the eigenvalue deficit $\epsilon$ quantifies the overlap.
\end{remark}

\subsection{Sign structure of the correction}

The geometric correction $\delta\Phi$ (Definition~\ref{def:delta-phi}) is strictly positive, since each adjacent overlap satisfies $C_{d,d+1}<1$ by Cauchy--Schwarz. Applied multiplicatively as $Q = Q_0\cdot\exp(\delta\Phi)$ (Corollary~\ref{cor:exp-resum}), it always increases the predicted value. This produces the observed sign structure across all cascade predictions:

\begin{itemize}
\item \emph{Descent-dependent quantities} (those traversing multiple cascade layers) have negative leading-order deviations, because the independent-step approximation underestimates the propagated geometric content. The correction $\exp(+\delta\Phi)$ closes the gap.
\item \emph{Geometric quantities} (single-level ratios, topological constants) have positive leading-order deviations, because they do not traverse multiple layers and receive no inter-layer correction. Their deviations reflect the subleading structure of the Gamma function at a single level, which is a separate (and smaller) effect.
\end{itemize}

The clean sign separation---negative for descent, positive for geometric---is not an empirical observation. It is a consequence of $\delta\Phi > 0$, which is itself a consequence of $C_{ij} < 1$ for $i\neq j$ (Cauchy--Schwarz on non-collinear cascade integrands).

\subsection{The correction as compactification residual}

Theorem~\ref{thm:reff} establishes that each compactification step retains a residual $R_{\mathrm{eff}}(d) = 1/\sqrt{d+3}$. The overlap deficit (Theorem~\ref{thm:overlap}) scales as $1-C_{d,d+1}^2 \sim 1/(8d^2) \sim R_{\mathrm{eff}}^4/8$. The eigenvalue deficit $\epsilon$ is the cumulative effect of these per-step residuals over the path.

The independent-step approximation corresponds to the limit $R_{\mathrm{eff}} \to 0$ (complete compactification). The geometric correction $\delta\Phi$ (Definition~\ref{def:delta-phi}) is the leading-order cost of $R_{\mathrm{eff}} > 0$---the imperfect folding of the cascade's own geometry.

\subsection{Numerical evidence and scope}

Theorem~\ref{thm:first-order} and Corollary~\ref{cor:exp-resum} (Section~\ref{sec:first-order} above) close the inter-layer correction at first order. Six descent-dependent cascade observables, each with a structurally-fixed path, are gathered below to show how the derived $\exp(\delta\Phi)$ correction performs against observation:

\begin{center}
\renewcommand{\arraystretch}{1.3}
\begin{tabular}{llcccccc}
\toprule
Observable & Path & $n$ & $\sum(1-C^2)$ & Leading & Corrected & Observed & Dev \\
\midrule
$\alpha_s$ & $d\!=\!5$--$12$ & 8 & 0.01186 & 0.1159 & 0.1173 & 0.1179 & $-0.5\%$ \\
$m_\tau/m_\mu$ & $d\!=\!6$--$13$ & 8 & 0.00938 & 16.53 & 16.69 & 16.82 & $-0.8\%$ \\
$m_\mu/m_e$ & $d\!=\!14$--$21$ & 8 & 0.00272 & 206.50 & 207.06 & 206.77 & $+0.1\%$ \\
$m_\tau$~(MeV) & $d\!=\!5$--$12$ & 8 & 0.01186 & 1755 & 1776 & 1777 & $-0.1\%$ \\
$v$~(GeV) & $d\!=\!5$--$12$ & 8 & 0.01186 & 240.8 & 243.7 & 246.2 & $-1.0\%$ \\
$\Omega_m^{\rm Bott}$ & $d\!=\!5$--$217$ & 213 & 0.02108 & 0.3115 & 0.3181 & 0.315 & $+1.0\%$ \\
\bottomrule
\end{tabular}
\end{center}

\noindent The first-order correction reduces descent deviations from $\sim\!2\%$ to sub-$1\%$ for all tested quantities. The correction for $m_\mu/m_e$ is small ($0.27\%$) because its path lies at higher $d$ where the per-step deficit scales as $1/(8d^2)$.

\begin{remark}[Scope of Theorem~\ref{thm:first-order}; status of $\Omega_m^{\rm Bott}$]\label{rem:sp9-status}
Theorem~\ref{thm:first-order} is proved for a \emph{multiplicative
propagator} traversing a contiguous path---i.e., for an observable of
the form $Q = \prod_{k=0}^{n-1}(\text{step factor at }d_k)$. Five of the
six table entries ($\alpha_s$, $m_\tau/m_\mu$, $m_\mu/m_e$, $m_\tau$,
$v$) have this structure directly. The sixth entry,
$\Omega_m^{\rm Bott}$, is structurally different: by
Part~V's ``Bott partition matter fraction'' theorem,
\[
\Omega_m^{\rm Bott}
\;=\;
\frac
{\displaystyle\sum_{\substack{d=5\\ d\bmod 8\in\{5,6\}}}^{217}\Omega_{d-1}}
{\displaystyle\sum_{d=5}^{217}\Omega_{d-1}},
\]
a \emph{ratio of sums} over cascade layers rather than a single
multiplicative propagator. Its inclusion in the table under the
full-path correction $\sum_{d=5}^{217}(1-C^2_{d,d+1}) = 0.02108$
applies the theorem outside its proven scope:
\begin{itemize}
\item The rigorous per-term treatment would correct each
$\Omega_{d-1}$ in the sums by its own path-length-dependent factor,
$\Omega_{d-1}\mapsto\Omega_{d-1}\cdot
\exp\bigl(\sum_{k=4}^{d-1}(1-C^2_{k,k+1})\bigr)$, then re-form the
ratio. The sums are dominated by the early terms near $d\sim 5$--$7$
(where $\Omega_d$ peaks), which have short paths and small corrections
($\sim\exp(0.001)\approx 1.001$). The resulting correction to the
ratio is therefore much smaller than the full-path factor
$\exp(0.02108)$.
\item The observed numerics support this: the leading value
$\Omega_m^{\rm Bott,\;lead}=0.31150$ sits $-1.1\%$ below the
observed $0.315$; the full-path ``corrected'' value $0.3181$
overshoots by $+1.0\%$. A rigorous per-term correction would fall
between these and is consistent with the two-population systematic
to leading order. The table entry reports the uniform full-path
correction for comparability with the other rows; it should be read
as an upper bound on the Gram shift for this observable, not as its
derived value.
\end{itemize}
A formal extension of Theorem~\ref{thm:first-order} to ratio-of-sums
observables---deriving the effective correction for
$Q=\sum_i a_i/\sum_j b_j$ where each $a_i,b_j$ is a multiplicative
propagator---is Part~V Open Question~1 (and remains open: the
identification step from non-trivial-phase layers to matter is the
other half of that question). The $\Omega_m$ closure does not
\emph{require} the Gram correction: both the leading lapse identity
$\Omega_m=1/\pi=0.31831$ ($+1.1\%$, geometric population) and the
leading Bott value $0.31150$ ($-1.1\%$, descent-dependent population)
bracket observation to the cascade's stated precision.
\end{remark}

\subsection{What this section proves}

\begin{enumerate}
\item The Gram matrix of cascade layer integrands is a matrix of Beta function values (Theorem~\ref{thm:gram}).
\item Adjacent layers have a strictly positive overlap deficit $1 - C_{d,d+1}^2 > 0$, scaling as $1/(8d^2)$ (Theorem~\ref{thm:overlap}). The correlation has the exact closed form $C^2_{d,d+1} = R(2d+2)^2 / [R(2d+1)\,R(2d+3)]$ in cascade slicing ratios (Theorem~\ref{thm:gram-R}), equivalently $\log C^2_{d,d+1} = -\tfrac{1}{2}\Delta^2 \log\alpha\big|_{2d+2}$ (Corollary~\ref{cor:gram-laplacian}): the Gram correction is the discrete log-Laplacian of the cascade compliance function $\alpha(d) = R(d)^2/4$ (Definition~\ref{def:compliance-action}) at doubled argument. The same compliance admits a marginal Green's function identity $G(d^*) - G(d^*+1) = \alpha(d^*)$ (Proposition~\ref{prop:marginal-greens}) which is the linear source response used by Part~IVb's $\alpha(d^*)/\chi^k$ precision closures.
\item The correlation matrix has one dominant eigenvalue $\lambda_1 \approx n(1 - \epsilon)$ with $\epsilon > 0$ (Theorem~\ref{thm:eigenvalue}).
\item The eigenvalue deficit is derived from first-order perturbation theory as a sum of Beta function deficits (Theorem~\ref{thm:pert}).
\item The sign of the correction (positive) explains the clean sign separation between descent and geometric predictions.
\item The first-order multiplicative correction $\delta Q / Q_0 = \sum(1-C^2_{d,d+1})$ (Theorem~\ref{thm:first-order}), uniquely promoted to the exponential form $Q = Q_0 \cdot \exp(\sum(1-C^2))$ by multiplicative composition under path-splitting (Corollary~\ref{cor:exp-resum}). Both ingredients are derived; no free parameter enters.
\item Every quantity in this section is a Gamma function identity. No physics enters.
\end{enumerate}

\subsection{Open questions}\label{sec:open}

\begin{enumerate}
\item \emph{Analytic formula for $\epsilon(n, d_0)$.} Partially closed. At $n = 2$ the dominant eigenvalue admits the exact closed form $\lambda_1 = 1 + R(2d_0+2)/\sqrt{R(2d_0+1)R(2d_0+3)}$ (Corollary~\ref{cor:lambda1-n2}). For $n \geq 3$ the power-trace bound family $U_k = (\mathrm{tr}\,C^k)^{1/k}$ (Corollary~\ref{cor:lambda1-power-trace}) provides closed-form cascade-native approximations to $\lambda_1$ converging at geometric rate $(\lambda_2/\lambda_1)^k$, reaching machine precision at $k = 6$ for cascade-physics paths. What remains genuinely open is whether a single-equation radical or special-function closed form (rather than a convergent family of upper bounds) exists at $n \geq 3$.
\item \emph{Higher-order corrections.} The subdominant eigenvalues $\lambda_3, \lambda_4, \ldots$ contribute at third order and beyond. For the $d = 5, \ldots, 12$ path, $\lambda_2/n = 0.81\%$ and $\lambda_3/n = 0.008\%$: the third-order correction is negligible. For longer paths, it may not be.
\item \emph{Application to mixed quantities.} The first-order correction (Theorem~\ref{thm:first-order} and Corollary~\ref{cor:exp-resum}) applies cleanly to quantities with a single unambiguous multiplicative cascade path. Quantities expressible as ratios of sums over cascade layers ($\Omega_m^{\rm Bott}$ in Part~V) require a per-term treatment articulated in Remark~\ref{rem:sp9-status} but not formally derived. A formal extension of the first-order theorem to such ratio-of-sums observables is open.
\end{enumerate}

\begin{remark}[Status of the correction]
The first-order correction is derived from cascade-native first principles. Theorem~\ref{thm:gram-R} and Corollary~\ref{cor:gram-laplacian} establish that the per-step Gram correction is the discrete log-Laplacian of the cascade scalar action's compliance function $\alpha(d)=R(d)^2/4$, evaluated at the doubled cascade argument. Theorem~\ref{thm:first-order} sums this to the path-cumulative shift $\delta\Phi=\sum(1-C^2_{d,d+1})$, and Corollary~\ref{cor:exp-resum} promotes it to the unique closed form $Q = Q_0\cdot\exp(\delta\Phi)$ respecting multiplicative composition under path-splitting. Both single-source $\alpha(d^*)/\chi^k$ family of Part~IVb and the path-distributed Gram correction here therefore derive from the same compliance $\alpha(d)=R(d)^2/4$: source response and discrete log-Laplacian, respectively, of one structural object.
\end{remark}

\begin{remark}[Superseded for the eight precision observables]\label{rem:superseded-by-family}
For eight Standard Model precision observables---$\alpha_s(M_Z)$, $m_\tau/m_\mu$, $m_\tau$ absolute, $\ell_A$, $\sin^2\theta_W$, $\Omega_m$, $\theta_C$, and $b/s$---the Gram-matrix framework above is superseded by the sharper structural identity of Part~IVb. Both deviations in these observables are closed \emph{at experimental precision} by cascade potential shifts of the form $\delta\Phi = \alpha(d^*)/\chi^k$, where $d^*$ is a distinguished dimension from this paper's four-dimension tower (Theorem~\ref{thm:tower}) and $\chi=\chi(S^{2n})=2$ is the Euler characteristic:
\begin{itemize}
\item $\alpha(14)/\chi$ closes $\alpha_s$ and $m_\tau/m_\mu$ together;
\item $\alpha(19)/\chi$ closes the $\tau$ absolute mass and $\ell_A$ at the phase-transition threshold $d_1=19$;
\item $\alpha(5)/\chi^3$ closes $\sin^2\theta_W$ and $\Omega_m$ at the volume maximum $d_V=5$;
\item $\alpha(7)/\chi^2$ closes $\theta_C$ and $\alpha(7)/\chi^4$ closes $b/s$, both sourced at the area maximum $d_0=7$.
\end{itemize}
For these eight observables the cascade closes at experimental precision via a single-source mechanism that does not invoke the Gram framework at all. The first-order Gram correction here remains the appropriate cascade-native account for the remaining descent-dependent observables ($m_\mu/m_e$ to $0.13\%$, $v$ before its non-perturbative $\exp(-\pi/\alpha(5))$ closure, etc.), where it reduces deviations from $\sim 2\%$ to sub-$1\%$.
\end{remark}

% ============================================================
\section{Summary}\label{sec:summary}
% ============================================================

\begin{center}
\begin{tabular}{cll}
\toprule
Step & Statement & Source \\
\midrule
1 & $V_\infty = \Omega_\infty = 0$ & Definition \\
2 & Orthogonal $\Rightarrow\sqrt{\pi}$, irreducible & Thms~\ref{thm:sqrtpi},~\ref{thm:irred} \\
3 & $\Omega_{d-1}$ unique independent quantity & Cor~\ref{cor:primary} \\
4 & $d_V = 5$: volume maximum & argmax $V_d$ \\
5 & $d_0 = 7$: decay-rate zero $=$ max-boundary ball & Thm~\ref{thm:dual} \\
6 & $c_1 = \tfrac{1}{2}\ln\pi$; unique across four classes & Thms~\ref{thm:opA},~\ref{thm:c1unique} \\
7 & $c_2 = \sqrt{\pi}$: second canonical value & Thm~\ref{thm:twovalues} \\
8 & Orbit terminates (first-order $\Rightarrow$ two forms) & Thm~\ref{thm:terminate} \\
9 & $d_1 = 19$,\; $d_2 = 217$ & Thm~\ref{thm:thresholds} \\
10 & Four distinguished dimensions; no fifth & Thm~\ref{thm:tower} \\
11 & Regime partition: Growth/Decay/ExpDecay/Oblivion & \S\ref{sec:regimes} \\
12 & Equilibria structurally stable; thresholds sensitive & Thms~\ref{thm:stable},~\ref{thm:sensitive} \\
13 & Content--scale separation from stability & Thm~\ref{thm:cs} \\
14 & Content: product $\Omega_{19}\times\Omega_{217}$ & Thm~\ref{thm:content} \\
15 & Scale: ratio $\Omega_5/\Omega_7 = 3/\pi$; exponent~2 & Thm~\ref{thm:twoscales}; Lem~\ref{lem:bilinear} \\
16 & $I = 9\,\Omega_{19}\,\Omega_{217}/\pi^2$; unique bilinear form & Thm~\ref{thm:invariant} \\
17 & $I = \sup_{\text{labelings}}$ over boundary pairs; argmax $(7,19,217)$ & Thm~\ref{thm:variational} \\
18 & $I = 1.0989\times 10^{-120}$ & Computation \\
19 & $120\approx\pi e^{2\sqrt{\pi}}(2\sqrt{\pi}-1)/\ln 10$ & Thm~\ref{thm:hierarchy} \\
\bottomrule
\end{tabular}
\end{center}

Every step is a theorem about the Gamma function. The derivation has zero
free parameters and no non-deductive steps.

% ============================================================
\section{What This Paper Does Not Do}\label{sec:what-this-paper-does-not-do}
% ============================================================

\begin{itemize}
\item Explain the agreement with
$\rho_\Lambda/M^4_{\mathrm{Pl}}$.
\item Introduce an observer, a hypothesis, or a physical interpretation.
\item Use any result from physics.
\end{itemize}

The paper is a theorem about the Gamma function.

% ============================================================
\section*{Appendix: Numerical Verification}
% ============================================================

$\sqrt{\pi} = 1.7724538509$\;;\quad
$\tfrac{1}{2}\ln\pi = 0.5723649429$

$\Omega_5 = \pi^3 = 31.006$\;;\quad
$\Omega_6 = 16\pi^3/15 = 33.073$\;;\quad
$\Omega_7 = \pi^4/3 = 32.470$

Discrete argmax at $d = 6$:\;$\Omega_6 > \Omega_7 > \Omega_5$\;$\checkmark$

$\Omega_5/\Omega_7 = 3/\pi = 0.95493$

$p(6) = -0.0208 < 0 < 0.0557 = p(7)$\;$\checkmark$\quad
(decay-rate zero between 6 and 7)

$\Omega_{19} = 0.51614$\;;\quad
$\Omega_{217} = 2.33490\times 10^{-120}$

$p(19) = 0.55351 < 0.57236 < p(20) = 0.57914$\;$\checkmark$
\quad(continuous crossing $d_1^\star\approx 19.7308$)

$p(217) = 1.77101 < 1.77245 < p(218) = 1.77331$\;$\checkmark$
\quad(continuous crossing $d_2^\star\approx 217.6267$)

$I_0 = 0.51614\times 2.33490\times 10^{-120}
= 1.2051\times 10^{-120}$\;$\checkmark$

$9/\pi^2 = 0.91189$\;$\checkmark$

$I = 0.91189\times 1.2051\times 10^{-120}
= 1.0989\times 10^{-120}$\;$\checkmark$

$\pi e^{2\sqrt{\pi}}(2\sqrt{\pi}-1)/\ln 10 = 120.27$\;$\checkmark$

$R_{\mathrm{eff}}(5) = 1/\sqrt{8} = 0.354$\;;\quad
$R_{\mathrm{eff}}(19) = 1/\sqrt{22} = 0.213$\;;\quad
$R_{\mathrm{eff}}(217) = 1/\sqrt{220} = 0.067$

$\Omega_4/V_5 = 5$\;$\checkmark$\;;\quad
$\Omega_6/V_7 = 7$\;$\checkmark$

\textbf{Structural stability verification:}

\begin{center}
\begin{tabular}{cccc}
\toprule
$K$ & $d_0(K)$ & $d_1(K)$ & $d_2(K)$ \\
\midrule
$1.50$ & 5 & 11 & 91 \\
$\sqrt{\pi} = 1.77$ & 7 & 19 & 217 \\
$2.00$ & 8 & 32 & 437 \\
$2.50$ & 13 & 79 & 1856 \\
\bottomrule
\end{tabular}
\end{center}

Equilibria exist for all $K$; thresholds shift by orders of
magnitude.\;$\checkmark$

\end{document}
